\documentclass[twocolumn,11pt,a4paper]{article}

\usepackage[utf8]{inputenc}
\usepackage[T1]{fontenc}
\usepackage{lmodern}
\usepackage[margin=1in]{geometry}
\usepackage{amsmath,amssymb,amsthm,mathtools}
\usepackage{thmtools}
\usepackage{enumitem}
\usepackage{microtype}
\usepackage{booktabs}
\usepackage{graphicx}
\usepackage{caption}
\usepackage[colorlinks=true,
            linkcolor=blue!50!black,
            citecolor=green!50!black,
            urlcolor=blue!60!black]{hyperref}
\usepackage{cleveref}
\usepackage{xcolor}

%------------------------------------------------------------------
%  Theorem environments (thmtools)
%------------------------------------------------------------------
\declaretheoremstyle[
  spaceabove=6pt, spacebelow=6pt,
  headfont=\bfseries, notefont=\mdseries, notebraces={(}{)},
  bodyfont=\itshape, postheadspace=1em
]{thmstyle}

\declaretheoremstyle[
  spaceabove=6pt, spacebelow=6pt,
  headfont=\bfseries, notefont=\mdseries, notebraces={(}{)},
  bodyfont=\normalfont, postheadspace=1em
]{defstyle}

\declaretheorem[style=thmstyle, name=Theorem,     numberwithin=section]{theorem}
\declaretheorem[style=thmstyle, name=Proposition, sibling=theorem]{proposition}
\declaretheorem[style=thmstyle, name=Lemma,       sibling=theorem]{lemma}
\declaretheorem[style=thmstyle, name=Corollary,   sibling=theorem]{corollary}
\declaretheorem[style=defstyle, name=Definition,  sibling=theorem]{definition}
\declaretheorem[style=defstyle, name=Axiom,       sibling=theorem]{axiom}
\declaretheorem[style=defstyle, name=Remark,      sibling=theorem]{remark}
\declaretheorem[style=defstyle, name=Example,     sibling=theorem]{example}
\declaretheorem[style=defstyle, name=Principle,   sibling=theorem]{principle}

% Named unnumbered environment for economic interpretations
\newenvironment{econinterp}[1]{%
  \medskip\noindent\textbf{Economic Interpretation #1.}\normalfont\itshape\space
}{%
  \medskip
}

%------------------------------------------------------------------
%  Macros
%------------------------------------------------------------------
\newcommand{\R}{\mathbb{R}}
\newcommand{\N}{\mathbb{N}}
\newcommand{\Z}{\mathbb{Z}}
\newcommand{\X}{\mathcal{X}}
\newcommand{\Y}{\mathcal{Y}}
\newcommand{\K}{\mathcal{K}}
\newcommand{\Cell}{\mathcal{C}}
\newcommand{\Method}{\mathfrak{M}}
\newcommand{\Goal}{\mathfrak{G}}
\newcommand{\receiver}{\mathsf{R}}
\newcommand{\agent}{\mathsf{A}}
\newcommand{\Sfloor}{S_{\flat}}
\newcommand{\eps}{\varepsilon}
\newcommand{\norm}[1]{\lVert#1\rVert}
\newcommand{\abs}[1]{\lvert#1\rvert}
\newcommand{\diam}{\mathrm{diam}}
\newcommand{\dom}{\mathrm{dom}}
\newcommand{\id}{\mathrm{id}}
\newcommand{\dd}{\,\mathrm{d}}
\newcommand{\Action}{\mathrm{Act}}
\newcommand{\Reach}{\mathrm{Reach}}
\newcommand{\Compose}{\diamond}

%------------------------------------------------------------------
%  Front matter
%------------------------------------------------------------------
\title{\textbf{A Mathematical Theory of Economic Agents:\\
Receivers, Floors, and the Algebra of Bounded Inquiry}}

\author{%
  Kundai Farai Sachikonye\\[4pt]
  \small Technical University of Munich\\
  \small AIMe Registry for Artificial Intelligence\\
  \small Maximus-von-Imhof-Forum~3, 85354 Freising, Germany\\
  \small \href{mailto:kundai.sachikonye@wzw.tum.de}{kundai.sachikonye@wzw.tum.de}
}

\date{}

%==================================================================
\begin{document}
%==================================================================
\maketitle

\begin{abstract}
We develop a mathematical theory of economic agents from a single primitive: the
\emph{receiver}, a four-tuple $(K,D,\Pi,\beta)$ consisting of a knowledge
framework, a decoder, a candidate-projection, and a strictly positive noise
floor.  From this primitive we derive, as theorems rather than assumptions, the
following: (i) the Floor Theorem — the irreducible lower bound $\beta>0$ on any
receiver's semantic distance to a target; (ii) the Cell-Truth Theorem — all
states within the same action-cell are $S$-indistinguishable; (iii) Mode
Non-Privilege — the priority of fast pre-decoder layers is a theorem of
optimality, not a heuristic assumption; (iv) a Banach fixed-point characterisation
of every methodology's floor; (v) the Agent Triple $(R,M,G)$ and its composite
floor; (vi) bounded rationality as a mathematical theorem; (vii) eleven classical
prerequisites for point-meaning each individually entail $\beta=0$, which the
Floor Theorem forbids; (viii) the \emph{G\"odelian residue} — private information
irreducible below $\beta$ — as a structural property of any bounded receiver; and
(ix) Arrow's impossibility, the Grossman--Stiglitz paradox, and the
exploration--exploitation incompatibility as corollaries.  All results are proved
in full.  Thirty-five numerical experiments verify every quantitative prediction
to machine precision ($\leq 10^{-14}$).
\end{abstract}

{\small\tableofcontents}

%==================================================================
\section{Introduction}
\label{sec:intro}
%==================================================================

\subsection{The agent problem in economics}
\label{sec:intro:problem}

The canonical economic agent is specified by a utility function
$U\colon X\to\R$ defined over a set of outcomes $X$, together with a
complete and transitive preference ordering $\succeq$ over $X$ induced
by the relation $x\succeq x'\iff U(x)\geq U(x')$.  Behaviour is then
the maximisation of $U$ subject to a budget or technological constraint.
This architecture, developed axiomatically by von Neumann and Morgenstern
\cite{vnm1944} and extended by Savage \cite{savage1954} to subjective
probability, underlies virtually all of positive economics.

Three foundational difficulties attend this specification.

\textbf{Preferences are assumed, not derived.}  The ordering $\succeq$ is
taken as a primitive datum of the model.  No mechanism is given for
\emph{generating} preferences from an underlying representation system.
When preferences conflict across contexts or agents, the standard apparatus
provides no principled means of resolution.

\textbf{Bounded rationality is a behavioral patch.}  Simon
\cite{simon1955,simon1957} documented that real agents do not globally
maximise but rather \emph{satisfice}: they search sequentially through
alternatives and halt at the first option that meets an aspiration level.
The standard response has been to append ``bounds'' to the rational-choice
model — computational limits, search costs, attention constraints — as
auxiliary hypotheses.  These bounds are behaviorally motivated; they are
not derived from the agent's representational structure.

\textbf{Substrate neutrality is lacking.}  A human consumer, a firm, and an
algorithmic trading system receive the same formal apparatus: a utility
function and a constraint set.  Nothing in the standard theory explains why
entities with radically different computational substrates and representational
structures should share an identical mathematical description.

\textbf{Point-valued utility requires infinite precision.}  The function
$U\colon X\to\R$ assigns a unique real number to every outcome state.
This presupposes that every state is uniquely identifiable — a form of
infinite representational precision that we show is impossible for any
bounded agent.

\medskip

The present paper addresses all four difficulties by deriving the agent
from a simpler primitive: the \emph{receiver}.  We show that bounded
rationality, selective rationality, the impossibility of point-values, the
irreducibility of private information, and the structural limits of value
theory are all theorems that follow from the receiver's definition.  They
are not modifications of the standard model; they replace it.

\subsection{Main contributions}
\label{sec:intro:contributions}

The contributions of this paper are as follows.

\begin{enumerate}[label=(\roman*),leftmargin=*]
\item \textbf{Receiver primitive and Floor Theorem} (\Cref{sec:receivers,sec:floor}).
  We introduce the receiver $(K,D,\Pi,\beta)$ and prove that
  $\Sfloor(R)=\beta>0$ is the irreducible lower bound on the
  $S$-functional for any bounded receiver.

\item \textbf{Cell-Truth Theorem} (\Cref{sec:celltruth}).  All states
  within the same action-cell attain $S=\beta$ and are therefore
  $S$-indistinguishable.  Representational invariance across isometric
  encodings follows as a corollary.

\item \textbf{Mode Non-Privilege and Selective Rationality}
  (\Cref{sec:layered}).  For a layered receiver, the pre-decoder layer
  (System~1) has strict priority whenever it suffices.  This is a theorem
  of optimality, not a behavioral assumption.

\item \textbf{Methodology Banach Fixed Point} (\Cref{sec:methodology}).
  Every methodology $M=(T,\kappa,\sigma)$ has a unique attractor
  $\Sfloor(M)=\sigma\kappa/(1-\kappa)>0$ that is the irreducible
  methodological floor.

\item \textbf{Agent Triple} (\Cref{sec:agent}).  The agent $A=(R,M,G)$ has
  floor $\Sfloor(A)=\beta\cdot\sigma\kappa/((1-\kappa)\Sigma)$.  The
  construction--action exclusion principle follows.

\item \textbf{Bounded Rationality as Theorem} (\Cref{sec:floor}).
  No bounded receiver can achieve $S=0$.  Simon's satisficing is optimal by
  \Cref{thm:selectiverationality}.

\item \textbf{Eleven-fold collapse to $\beta=0$} (\Cref{sec:cellvalue}).
  Eleven classical prerequisites for point-meaning each individually require
  $\beta=0$, which the Floor Theorem forbids.

\item \textbf{G\"odelian residue $=$ private information $=$ $\beta$}
  (\Cref{sec:godelian}).  The residual content of any solution that the
  agent cannot communicate is exactly the floor $\beta$.

\item \textbf{Thirty-five numerical experiments} (\Cref{sec:numerical}).
  All quantitative predictions are verified to relative error
  $\leq 10^{-14}$.
\end{enumerate}

\subsection{Organisation}
\label{sec:intro:org}

\Cref{sec:outcomes} introduces outcome spaces and action-cells.
\Cref{sec:receivers} defines the receiver and the $S$-functional.
\Cref{sec:floor} proves the Floor Theorem and establishes bounded
rationality.
\Cref{sec:celltruth} proves Cell-Truth and Representational Invariance.
\Cref{sec:layered} treats layered receivers and selective rationality.
\Cref{sec:methodology} develops the methodology Banach fixed point.
\Cref{sec:agent} defines the agent triple and the uncertainty principle.
\Cref{sec:cellvalue} proves the eleven-fold collapse and the cell-value
theorems.
\Cref{sec:godelian} treats the G\"odelian residue and private information.
\Cref{sec:impossibilities} reconsiders classical impossibility results.
\Cref{sec:numerical} presents the numerical validation.
\Cref{sec:litconnections} situates the theory in the existing literature.
\Cref{sec:conclusion} concludes.

%==================================================================
\section{Outcome Spaces and Action-Cells}
\label{sec:outcomes}
%==================================================================

\begin{definition}[Outcome space]
\label{def:outcomespace}
An \emph{outcome space} is a triple $(X,d,\Sigma)$ where
\begin{enumerate}[label=(\alph*),leftmargin=*]
  \item $X$ is a nonempty set of \emph{outcome states};
  \item $d\colon X\times X\to[0,\Sigma]$ is a \emph{bounded
    pseudo-metric}: $d(x,x)=0$ for all $x$; $d(x,x')=d(x',x)$; and
    $d(x,x'')\leq d(x,x')+d(x',x'')$ for all $x,x',x''\in X$;
  \item $\Sigma\in(0,\infty)$ is the \emph{maximal semantic distance}.
\end{enumerate}
The canonical normalisation is $\Sigma=100$; no result in this paper
depends on the specific value.
\end{definition}

\begin{definition}[Action map]
\label{def:actionmap}
An \emph{action map} is a measurable surjection
$\Action\colon X\to Y$, where $(Y,d_Y)$ is a metric space called the
\emph{action space}.  For any outcome state $x\in X$, the element
$\Action(x)\in Y$ is the \emph{practical-response class} of $x$.
\end{definition}

\begin{definition}[Action-cell]
\label{def:cell}
For $y\in Y$ the \emph{action-cell} of $y$ is the preimage
\[
  C_y \;:=\; \Action^{-1}(\{y\}) \;\subseteq\; X.
\]
Each action-cell is equipped with three numerical invariants:
\begin{enumerate}[label=(\alph*),leftmargin=*]
  \item \emph{Tolerance}: $\tau(C_y):=\sup_{x\in X}\inf_{c\in C_y}
    d(x,c)>0$.
  \item \emph{Cell distance}: $d(x,C_y):=\inf_{c\in C_y}d(x,c)$.
  \item \emph{Diameter}: $\diam(C_y):=\sup_{c,c'\in C_y}d(c,c')<\Sigma$.
\end{enumerate}
A \emph{generic action-cell} is any nonempty $C\subseteq X$ with
$\tau(C)>0$ and $\diam(C)<\Sigma$.
\end{definition}

\begin{proposition}[Metric properties of cell distance]
\label{prop:celldistmetric}
For any generic action-cell $C\subseteq X$, the function
$x\mapsto d(x,C)$ is $1$-Lipschitz: for all $x,x'\in X$,
\[
  \abs{d(x,C)-d(x',C)}\;\leq\;d(x,x').
\]
Moreover $d(x,C)=0$ for all $x\in C$, and $d(x,C)\leq\Sigma$ globally.
\end{proposition}

\begin{proof}
Let $x,x'\in X$ and $\eps>0$.  Choose $c\in C$ such that
$d(x,c)\leq d(x,C)+\eps$.  Then
\[
  d(x',C)\;\leq\;d(x',c)\;\leq\;d(x',x)+d(x,c)\;\leq\;d(x',x)+d(x,C)+\eps.
\]
Since $\eps>0$ was arbitrary, $d(x',C)\leq d(x,x')+d(x,C)$.  By
symmetry $d(x,C)\leq d(x,x')+d(x',C)$, giving the Lipschitz bound.
For $x\in C$ take $c=x$ in the infimum: $d(x,C)\leq d(x,x)=0$, so
$d(x,C)=0$.  The upper bound $d(x,C)\leq\Sigma$ follows because $d$
takes values in $[0,\Sigma]$.
\end{proof}

\begin{remark}[Operational granularity]
\label{rem:granularity}
The condition $\tau(C)>0$ expresses that the cell has positive
\emph{operational granularity}: there exist states not in $C$ whose
distance to $C$ is bounded away from zero by a uniform constant.  A cell
of zero tolerance is a singleton.  As we prove in \Cref{sec:floor}
(\Cref{cor:noperfectprecision}), singletons are unreachable by any bounded
receiver.
\end{remark}

\begin{econinterp}{2.6 (Goods markets)}
In a goods market, $X$ is the set of possible transaction states
(prices, quantities, qualities, timing), $Y$ is the set of distinct
practical transaction outcomes, and $\Action$ maps each state to its
transaction class.  The action-cell $C_y$ is the set of states that
produce the same practical transaction outcome $y$ — the ``price band''
within which exchange is operationally indistinguishable.  The diameter
$\diam(C_y)$ is the width of the price band; the tolerance $\tau(C_y)$
is the minimum clearing distance between bands.
\end{econinterp}

%==================================================================
\section{Receivers and the $S$-Functional}
\label{sec:receivers}
%==================================================================

\begin{definition}[Receiver]
\label{def:receiver}
A \emph{receiver} on $(X,d,\Sigma)$ is a four-tuple
$R=(K,D,\Pi,\beta)$ where:
\begin{enumerate}[label=(\alph*),leftmargin=*]
  \item $K$ is a nonempty set called the \emph{knowledge framework};
  \item $D\colon X\to K$ is the \emph{decoder}, a measurable map;
  \item $\Pi\colon K\to 2^X\setminus\{\emptyset\}$ is the
    \emph{candidate-projection}, satisfying the \emph{consistency
    condition}: for all $x\in X$,
    \[
      D(x)\;\in\;D\!\left(\Pi(D(x))\right),
    \]
    meaning the decoded representation of any candidate lies in the
    same decoded-image class as the original;
  \item $\beta\in(0,\Sigma]$ is the \emph{noise floor},
    \[
      \beta\;:=\;\sup_{x\in X}\;\inf_{x'\in\Pi(D(x))} d(x,x')\;>\;0.
    \]
\end{enumerate}
\end{definition}

\begin{figure*}[!htbp]
    \centering
    \includegraphics[width=\textwidth]{panels/paper1_panel_1.png}
    \caption{\textbf{S-Functional Geometry: Floor Surface, Spatial Field, Family of Floors, and Cell-Truth Attainment.}
    (\textbf{A})~Three-dimensional surface of the S-functional $S(\beta, d) = \max(\beta,\, d(x, C))$ over a grid of receiver floors $\beta \in [0.1, 4.0]$ and distances $d(x,C) \in [0, 6]$.  The ridge along the diagonal $\beta = d$ separates the floor-dominated regime ($d < \beta$, where $S = \beta$ regardless of proximity to the cell) from the distance-dominated regime ($d > \beta$, where $S = d$). 
    (\textbf{B})~Spatial heatmap of $S(\mathcal{R}, x; C)$ for a ball cell $C = B(\mathbf{0}, r)$ with $r = 1.2$ and floor $\beta = 0.6$ in $\mathbb{R}^2$.  The white circle marks the cell boundary $\partial C$.  Inside $C$ the field is constant at $\beta$ (the Cell-Truth property: $x \in C \Rightarrow S = \beta$), while outside it grows monotonically with $\|x\|$, saturating to $\|x\| - r + \beta$ for large $\|x\|$. 
    (\textbf{C})~S-functional versus distance $d(x, C)$ for five floors $\beta \in \{0.3, 0.8, 1.5, 2.5, 3.8\}$.  Each curve is piecewise linear with a kink at $d = \beta$: when $d < \beta$ the curve is horizontal at height $\beta$ so the floor dominates, and when $d > \beta$ the curve rises with unit slope so distance dominates. 
    (\textbf{D})~Cell-Truth attainment: S-functional values at $N = 60$ randomly drawn in-cell states for 60 receivers with floors $\beta \in [0.2, 4.0]$, plotted against $\beta$.  Every point lies exactly on the diagonal $S = \beta$, confirming that the Cell-Truth property $S(\mathcal{R}, x; C) = \beta$ for all $x \in C$ holds with zero error across all floor values and all receiver configurations tested.  }
    \label{fig:agent_panel_1}
\end{figure*}

\begin{definition}[$S$-functional]
\label{def:sfunctional}
For receiver $R=(K,D,\Pi,\beta)$, outcome state $x\in X$, and generic
action-cell $C\subseteq X$, the \emph{$S$-functional} is
\[
  S(R,x;C)\;:=\;\inf_{x'\in\Pi(D(x))} d(x',C)\;+\;\beta.
\]
When the target cell $C=C_{\Action(x)}$ is clear from context we write
$S(R,x):=S(R,x;C_{\Action(x)})$.
\end{definition}

\begin{proposition}[Basic properties of $S$]
\label{prop:sproperties}
For any receiver $R=(K,D,\Pi,\beta)$ and generic action-cell $C$:
\begin{enumerate}[label=(\roman*),leftmargin=*]
  \item \emph{Floor:} $S(R,x;C)\geq\beta$ for all $x\in X$.
  \item \emph{Upper bound:} $S(R,x;C)\leq d(x,C)+\beta$.
  \item \emph{Monotonicity:} $C\subseteq C'$ implies
    $S(R,x;C')\leq S(R,x;C)$.
  \item \emph{Upper semicontinuity:} if $D$ and $\Pi$ are Borel
    measurable, then $x\mapsto S(R,x;C)$ is upper semicontinuous.
\end{enumerate}
\end{proposition}

\begin{proof}
\textit{(i):} Since $\beta>0$, we have
$S(R,x;C)=\inf_{x'\in\Pi(D(x))}d(x',C)+\beta\geq 0+\beta=\beta$.

\textit{(ii):} Taking $x'=x$ in the infimum over $\Pi(D(x))$ gives
$\inf_{x'\in\Pi(D(x))}d(x',C)\leq d(x,C)$, so
$S(R,x;C)\leq d(x,C)+\beta$.

\textit{(iii):} For $C\subseteq C'$ and any $x'\in\Pi(D(x))$,
$d(x',C')=\inf_{c'\in C'}d(x',c')\leq\inf_{c\in C}d(x',c)=d(x',C)$
since $C\subseteq C'$.  Taking the infimum over $\Pi(D(x))$ gives
$\inf_{x'}d(x',C')\leq\inf_{x'}d(x',C)$, hence
$S(R,x;C')\leq S(R,x;C)$.

\textit{(iv):} The map $x\mapsto\inf_{x'\in\Pi(D(x))}d(x',C)$ is the
infimum of a family of Borel functions indexed by the Borel image of
$\Pi\circ D$.  The infimum of an upper-semicontinuous family of
measurable functions is upper semicontinuous \cite[Thm.~7.14]{rudin1976}.
Adding the constant $\beta$ preserves upper semicontinuity.
\end{proof}

\begin{definition}[Receiver floor]
\label{def:receiverfloor}
The \emph{receiver floor} is $\Sfloor(R):=\beta$.
\end{definition}

\begin{econinterp}{3.5 (Perception and inference)}
The knowledge framework $K$ is the set of concepts and representations
available to the agent — its internal language for describing states.
The decoder $D$ maps observable states to internal representations: this
is \emph{perception}.  The candidate-projection $\Pi$ maps each
representation back to the set of outcome states compatible with that
representation: this is \emph{inference}.  The floor $\beta$ is the
minimum remaining uncertainty after both perception and inference:
even when the agent correctly recognises a state, there remain states at
distance $\beta$ that are indistinguishable under its decoder.  This is
not a modelling error; it is a structural property of bounded cognition.
\end{econinterp}

%==================================================================
\section{The Floor Theorem and Bounded Rationality}
\label{sec:floor}
%==================================================================

\begin{theorem}[Floor Theorem]
\label{thm:floor}
Let $R=(K,D,\Pi,\beta)$ be a receiver on $(X,d,\Sigma)$.  Then
$\Sfloor(R)=\beta>0$ is the irreducible lower bound on $S(R,x;C)$
for every $x\in X$ and every generic action-cell $C$.  Moreover, the
bound is attained: if $x\in C$ and the decoder-projection chain is
faithful at $x$ (i.e.\ $d(x,C)=0$), then $S(R,x;C)=\beta$.
\end{theorem}

\begin{proof}
The lower bound $S(R,x;C)\geq\beta$ is \Cref{prop:sproperties}(i).

For attainment, suppose $x\in C$.  Then $d(x,C)=\inf_{c\in C}d(x,c)
\leq d(x,x)=0$, so $d(x,C)=0$.  By \Cref{prop:sproperties}(ii),
$S(R,x;C)\leq d(x,C)+\beta=0+\beta=\beta$.  Combined with the lower
bound, $S(R,x;C)=\beta$.
\end{proof}

\begin{corollary}[No perfect precision]
\label{cor:noperfectprecision}
No bounded receiver can achieve $S(R,x;C)=0$.  Any finite agent retains
residual semantic distance at least $\beta>0$ to any target cell.
\end{corollary}

\begin{proof}
By \Cref{def:receiver}, $\beta>0$ is part of the receiver's definition.
By \Cref{prop:sproperties}(i), $S(R,x;C)\geq\beta>0$.
\end{proof}

\begin{theorem}[Bounded rationality is a theorem]
\label{thm:boundedrationalitytheorem}
Let $A=(R,M,G)$ be an agent (see \Cref{def:agent}) with receiver $R$
having floor $\beta>0$.  Then for any generic action-cell $C$ with
$\tau(C)>0$, the irreducible residual $\Sfloor(R)=\beta>0$ cannot be
eliminated by any finite computation within the methodology $M$.
Specifically, for all finite iterates $t\geq 1$,
\[
  S_t(R,x;C)\;\geq\;\beta\;>\;0.
\]
\end{theorem}

\begin{proof}
The floor $\beta$ is a property of the receiver structure $(K,D,\Pi)$,
not of the methodology $M=(T,\kappa,\sigma)$.  The methodology produces
iterates that converge to the methodological floor $\Sfloor(M)$ (proved
in \Cref{thm:banach}), but the receiver floor $\beta$ is a lower bound
on $S(R,x;C)$ for every state and every cell by
\Cref{prop:sproperties}(i).  Any finite computation within $M$ produces
candidates $x'\in\Pi(D(x))$; the infimum of $d(x',C)$ over this set is
non-negative, and adding $\beta>0$ gives $S(R,x;C)\geq\beta>0$ for all
finite iterations.  The bound cannot be reduced by increasing the number
of iterations because $\beta$ is a structural constant of the receiver,
not a function of the iteration count.
\end{proof}

\begin{remark}[Simon's bounded rationality]
\label{rem:simon}
Simon \cite{simon1955,simon1957} introduced bounded rationality as a
behavioral modification of expected utility theory: agents satisfice
rather than maximise because of computational and informational
constraints.  \Cref{thm:boundedrationalitytheorem} establishes this as a
mathematical theorem rather than a behavioral hypothesis.  The floor
$\beta$ is not a constraint appended to the model; it is derived from the
receiver structure.  The theorem extends Simon's observation from a
behavioral description to a mathematical necessity.
\end{remark}

\begin{remark}[Relationship to expected utility]
\label{rem:eurel}
Von Neumann and Morgenstern \cite{vnm1944} and Savage \cite{savage1954}
develop utility theory for agents with complete, transitive preferences
over outcomes.  The Floor Theorem establishes that any receiver
satisfying their framework (when embedded in an outcome space) necessarily
has $\beta>0$ and cannot achieve the point-precision their framework
assumes.  The utility function $U\colon X\to\R$ presupposes that every
$x$ maps to a unique value — a form of \emph{point-meaning} (defined in
\Cref{def:pointmeaning}) that we prove in \Cref{sec:cellvalue} is
forbidden.
\end{remark}

%==================================================================
\section{Cell-Truth and Representational Invariance}
\label{sec:celltruth}
%==================================================================

\begin{theorem}[Cell-Truth]
\label{thm:celltruth}
Let $\Action\colon X\to Y$ be measurable and let $C_y=\Action^{-1}(\{y\})$
be an action-cell with $\tau(C_y)>0$.  For any receiver $R$ and any
$x_1,x_2\in C_y$,
\[
  S(R,x_1;C_y)\;=\;S(R,x_2;C_y)\;=\;\beta.
\]
\end{theorem}

\begin{proof}
For any $x_i\in C_y$ we have $d(x_i,C_y)=0$ because $x_i\in C_y$.
By \Cref{prop:sproperties}(ii), $S(R,x_i;C_y)\leq 0+\beta=\beta$.
By \Cref{prop:sproperties}(i), $S(R,x_i;C_y)\geq\beta$.
Hence $S(R,x_i;C_y)=\beta$ for $i=1,2$.
\end{proof}

\begin{remark}
\label{rem:celltruthimport}
Cell-Truth is the fundamental theorem for value theory.  All states
within the same action-cell are $S$-indistinguishable: they produce the
same value of the $S$-functional.  Two states that produce identical
practical actions are operationally equivalent, regardless of how
different they appear in the representation space $K$.
\end{remark}

\begin{theorem}[Representational Invariance]
\label{thm:repinvariance}
Let $\varphi\colon X\to X'$ be an isometric bijection between two outcome
spaces.  Define the \emph{transferred receiver}
$R'=(K,D\circ\varphi^{-1},\,\varphi\circ\Pi,\,\beta)$ on $X'$ and the
transferred cell $C'=\varphi(C)$.  Then for all $x\in X$,
\[
  S(R,x;C)\;=\;S(R',\varphi(x);C').
\]
\end{theorem}

\begin{proof}
Cell membership is preserved by the bijection: $\varphi(x)\in C'$ if
and only if $x\in C$.  The infima coincide because $\varphi$ is an
isometry: for any $x'\in\Pi(D(x))$,
\[
  d_{X'}\!\left(\varphi(x'),\varphi(C)\right)
  \;=\;\inf_{c\in C}d_{X'}\!\left(\varphi(x'),\varphi(c)\right)
  \;=\;\inf_{c\in C}d_X(x',c)
  \;=\;d_X(x',C).
\]
The candidate set of $R'$ at $\varphi(x)$ is
$\Pi'(D'(\varphi(x)))=\varphi(\Pi(D(\varphi^{-1}(\varphi(x)))))
=\varphi(\Pi(D(x)))$.  Therefore
\[
  \inf_{z'\in\Pi'(D'(\varphi(x)))}d_{X'}(z',C')
  \;=\;\inf_{x'\in\Pi(D(x))}d_X(x',C).
\]
Adding $\beta$ (which is unchanged) gives $S(R,x;C)=S(R',\varphi(x);C')$.
\end{proof}

\begin{corollary}[Three canonical encodings]
\label{cor:threencodings}
The following three encodings of practical interest are isometric
embeddings of outcome spaces; Representational Invariance therefore
applies to each, and the $S$-functional is preserved across encodings.
\begin{enumerate}[label=(\roman*),leftmargin=*]
  \item \emph{Oscillatory} ($L^2$ embedding): outcomes embedded as
    square-integrable wave-packets; the metric is the $L^2$ norm.
  \item \emph{Categorical} (small-category morphism length): outcomes as
    objects; the metric is the minimum morphism-sequence length
    \cite{maclane1998}.
  \item \emph{Partition} (integer-partition $\ell^1$-metric): outcomes as
    integer partitions; the metric is the $\ell^1$ distance between
    sorted parts.
\end{enumerate}
The choice among these encodings is computational, not epistemic.
\end{corollary}

\begin{proof}
Each encoding defines an isometric embedding $\varphi_i\colon X\to X_i$
into a metric space $X_i$.  By \Cref{thm:repinvariance}, $S$ is
preserved under each $\varphi_i$.  Composing two such isometries gives
another isometry, so the $S$-value is invariant across all three
encodings simultaneously.
\end{proof}

\begin{econinterp}{5.5 (Re-parameterisation)}
Representational Invariance means that the operational content of an
economic state does not depend on the units, language, or framework used
to describe it.  A price expressed in dollars, in euros, or in terms of a
commodity basket produces identical $S$-values if the encoding is
isometric.  The standard economics practice of re-parameterising models
is justified by this theorem, not by convention.
\end{econinterp}

\begin{figure*}[!htbp]
    \centering
    \includegraphics[width=\textwidth]{panels/paper1_panel_2.png}
    \caption{\textbf{Representational Invariance: Distance Field, Rotation Isometry, Floor Family at the Cell Boundary, and Translation Isometry.}
    (\textbf{A})~Three-dimensional surface of the cell-distance function $d(x_1, x_2; C) = \max(0, \|x\| - r)$ for $C = B(\mathbf{0}, 1)$ in $\mathbb{R}^2$.  The flat plateau at $d = 0$ corresponds to the interior of $C$, where all states are semantically indistinguishable under any receiver.  The surrounding cone rises with unit slope in the Euclidean norm.  The surface's rotational symmetry about the vertical axis reflects the fact that $d(\cdot, C)$ depends only on $\|x - c_0\|$, not on the angular coordinate: cell-distance is invariant under any isometry that fixes the cell.
    (\textbf{B})~Representational Invariance (rotation): scatter of $S(\mathcal{R}, x; C)$ against $S(\mathcal{R}, \phi(x); \phi(C))$ for $N = 400$ query points and a 45-degree rotation $\phi$.  All points fall on the diagonal $y = x$ to within floating-point precision ($\max|\Delta S| < 10^{-13}$), confirming that $S$ is preserved under isometric bijections. 
    (\textbf{C})~S-functional versus $\|x\|$ (Euclidean norm of $x$) for four floors $\beta \in \{0.3, 0.8, 1.5, 2.5\}$.  The vertical dotted line at $\|x\| = r = 1$ marks the cell boundary.  Inside the cell ($\|x\| < r$), $S = \beta$ is constant.  Outside ($\|x\| > r$), $S$ increases linearly with slope 1.  
    (\textbf{D})~Representational Invariance (translation): scatter of $S(\mathcal{R}, x; C)$ against $S(\mathcal{R}, x + c; C + c)$ for $N = 400$ query points and shift $c = (2.3, -1.7)$.  }
    \label{fig:agent_panel_2}
\end{figure*}



%==================================================================
\section{Layered Receivers and Selective Rationality}
\label{sec:layered}
%==================================================================

\begin{definition}[Layered receiver]
\label{def:layered}
A \emph{layered receiver} is an ordered tuple
$R_\bullet=(R_1,\ldots,R_n)$ of receivers on the same outcome space
$(X,d,\Sigma)$, with aggregated $S$-functional
\[
  S(R_\bullet,x;C)\;:=\;\min_{1\leq i\leq n}S(R_i,x;C).
\]
Layer $i$ is called:
\begin{itemize}[leftmargin=*]
  \item \emph{pre-decoder} if $D_i$ is constant (a reflex or heuristic
    layer);
  \item \emph{decoder-layer} if $D_i$ is non-constant and continuous (a
    deliberative layer);
  \item \emph{delegated} if $D_i$ factors through an external receiver
    (outsourced reasoning).
\end{itemize}
\end{definition}

\begin{proposition}[Layered floor]
\label{prop:layeredfloor}
For a layered receiver $R_\bullet=(R_1,\ldots,R_n)$,
\[
  \Sfloor(R_\bullet)\;=\;\min_{1\leq i\leq n}\beta_i.
\]
\end{proposition}

\begin{proof}
By \Cref{prop:sproperties}(i), $S(R_i,x;C)\geq\beta_i$ for each layer
$i$ and all $x,C$.  Therefore
$S(R_\bullet,x;C)=\min_i S(R_i,x;C)\geq\min_i\beta_i$ for all $x,C$,
giving the lower bound.  The lower bound is attained by the layer $i^*$
that achieves $\min_i\beta_i$: at any $x\in C$,
$S(R_{i^*},x;C)=\beta_{i^*}=\min_i\beta_i$ by \Cref{thm:floor}, so
$S(R_\bullet,x;C)\leq\beta_{i^*}=\min_i\beta_i$.
\end{proof}

\begin{theorem}[Mode Non-Privilege]
\label{thm:modenonprivilege}
Let $R_\bullet$ be a layered receiver and let $R_{i^*}$ be a
pre-decoder layer with $\beta_{i^*}<\tau(C)$.  If $R_{i^*}$ fires at
state $x$ — meaning $S(R_{i^*},x;C)\leq\tau(C)$ — then
\[
  S(R_\bullet,x;C)\;\leq\;\beta_{i^*}\;<\;\tau(C),
\]
independently of the $S$-values of all decoder layers.
\end{theorem}

\begin{proof}
Since $S(R_\bullet,x;C)=\min_i S(R_i,x;C)$, we have
$S(R_\bullet,x;C)\leq S(R_{i^*},x;C)$.  The pre-decoder layer $R_{i^*}$
has constant decoder $D_{i^*}$, so $\Pi(D_{i^*}(x))$ is the same for all
$x$ (determined entirely by the constant value of $D_{i^*}$).  Because
$\beta_{i^*}$ is the floor of $R_{i^*}$,
$S(R_{i^*},x;C)\geq\beta_{i^*}$.  By hypothesis $S(R_{i^*},x;C)\leq
\tau(C)$, so $S(R_\bullet,x;C)\leq S(R_{i^*},x;C)\leq\tau(C)$.
Moreover $S(R_\bullet,x;C)\geq\min_i\beta_i\leq\beta_{i^*}$ (since
$i^*$ may or may not be the minimising layer).  In either case,
$S(R_\bullet,x;C)\leq\beta_{i^*}<\tau(C)$, and the decoder layers do
not influence this bound.
\end{proof}

\begin{theorem}[Selective rationality as optimality]
\label{thm:selectiverationality}
Let $R_\bullet$ be a layered receiver with target cell $C$ and tolerance
$\tau(C)$.  Suppose each layer $i$ has associated cost $w_i>0$.  The
optimal response — minimising total cost subject to achieving
$S(R_\bullet,x;C)\leq\tau(C)$ — is to use the cheapest layer $i^*$
with $\beta_{i^*}<\tau(C)$.  Using any more expensive layer when $i^*$
already achieves the target strictly increases cost without reducing
$S(R_\bullet,x;C)$ below $\beta_{i^*}$.
\end{theorem}

\begin{proof}
Let $i^*=\arg\min_{i:\beta_i<\tau(C)} w_i$ be the cheapest layer that
can achieve $S<\tau(C)$.  By \Cref{thm:modenonprivilege}, when $R_{i^*}$
fires, $S(R_\bullet,x;C)\leq\beta_{i^*}$.  For any other layer $j\neq
i^*$ with $w_j>w_{i^*}$, activating $j$ in addition to $i^*$ incurs
cost $w_j+w_{i^*}>w_{i^*}$ but cannot reduce
$S(R_\bullet,x;C)=\min_i S(R_i,x;C)$ below $\beta_{i^*}$, because
$S(R_{i^*},x;C)\leq\beta_{i^*}$ already sets the minimum.  Hence the
cheapest-layer strategy achieves the same $S$-value at strictly lower
cost; it is uniquely optimal.
\end{proof}

\begin{figure*}[!htbp]
    \centering
    \includegraphics[width=\textwidth]{panels/paper1_panel_3.png}
    \caption{\textbf{Layered Receivers and Selective Rationality: Aggregate S Surface, Layer Floor Hierarchy, Pre-Decoder Activation, and Dual-Process Transition.}
    (\textbf{A})~Three-dimensional surface of the aggregate S-functional $S(\mathcal{R}_\bullet, x; C) = \min_i S(\mathcal{R}_i, x; C) = \min(\max(\beta_1, d),\, \max(\beta_2, d))$ over floors $(\beta_1, \beta_2) \in [0.2, 3.0]^2$ for a fixed outer-state distance $d = 1.8$.  The surface exhibits a ridge along the diagonal $\beta_1 = \beta_2$ and descends monotonically toward lower floor values in either coordinate: adding a layer with a smaller floor strictly reduces the aggregate S. 
    (\textbf{B})~Bar chart of individual layer floors versus aggregate floor across $N = 12$ random configurations, each with four layers (colours) and floors drawn uniformly from $[0.2, 4.0]$.  The black bars show the aggregate floor $\min_i \beta_i$.  In every configuration the black bar is at or below all coloured bars, confirming that the aggregate floor equals the minimum layer floor.  
    (\textbf{C})~Pre-decoder activation fraction (System~1 firing rate) versus cell tolerance $\tau(C)$ for a receiver with $\beta_{\mathrm{pre}} = 0.5$ and $N = 800$ query points drawn uniformly from $[-5, 5]^2$. 
    (\textbf{D})~Decoder activation fraction (System~2 firing rate) versus $\tau(C)$, the complementary series $1 - \Pr[S_{\mathrm{pre}} \leq \tau]$.  The fraction is monotone non-increasing: larger cells reduce the demand for deliberate decoding.}
    \label{fig:agent_panel_3}
\end{figure*}

\begin{remark}[Kahneman's dual-process taxonomy]
\label{rem:kahneman}
Kahneman \cite{kahneman2011,kahnemantversky1979} distinguishes System~1
(fast, automatic, associative) from System~2 (slow, deliberative,
analytical) processing.  \Cref{thm:selectiverationality} derives this
distinction: System~1 corresponds to pre-decoder layers (small $\beta$,
low cost, fires automatically); System~2 to decoder layers (potentially
smaller $\beta$, higher cost, fires only when pre-decoder layers fail to
reach the cell tolerance).  Mode Non-Privilege (\Cref{thm:modenonprivilege})
derives, rather than assumes, the priority of System~1 processing when it
is sufficient.  This is not a heuristic bias but the mathematically
optimal response of a layered receiver.
\end{remark}

\begin{remark}[Satisficing]
\label{rem:satisficing}
Simon's satisficing \cite{simon1956} — choosing the first option that
exceeds an aspiration level rather than maximising — is the behavioral
description of a layered receiver using the minimum-cost layer that
achieves $S<\tau(C)$.  The aspiration level is $\tau(C)$, and
satisficing is optimal by \Cref{thm:selectiverationality}.
\end{remark}

%==================================================================
\section{Methodology and the Banach Floor}
\label{sec:methodology}
%==================================================================

\begin{definition}[Methodology]
\label{def:methodology}
A \emph{methodology} on $(X,d,\Sigma)$ is a triple $M=(T,\kappa,\sigma)$
where:
\begin{enumerate}[label=(\alph*),leftmargin=*]
  \item $T\colon X\times[0,\Sigma]\to[0,\Sigma]$ is a per-iteration
    update map;
  \item $\kappa\in[0,1)$ is the \emph{contraction coefficient};
  \item $\sigma\in[0,\Sigma]$ is the \emph{per-iteration dispersion};
\end{enumerate}
satisfying, for all $x\in X$ and $s_1,s_2\in[0,\Sigma]$:
\[
  \abs{T(x,s_1)-T(x,s_2)}\;\leq\;\kappa\abs{s_1-s_2}
  \qquad\text{and}\qquad
  T(x,s)\;\geq\;\sigma\kappa.
\]
\end{definition}

\begin{definition}[Methodology iteration]
\label{def:methodologyiter}
The \emph{canonical linear form} of the update is $L(s)=\kappa s+\sigma\kappa$.
The \emph{iterated $S$-value} under $M$ from initial value $s_0\in[0,\Sigma]$
is the sequence $s_t=T(x,s_{t-1})$ for $t\geq 1$, approximated by the
linear iterates $s_t^L=L(s_{t-1}^L)$.
\end{definition}

\begin{theorem}[Banach fixed point for methodology]
\label{thm:banach}
The map $L(s)=\kappa s+\sigma\kappa$ on $[0,\Sigma]$ with $\kappa\in(0,1)$
is a contraction.  Its unique fixed point is
\[
  \Sfloor(M)\;:=\;\frac{\sigma\kappa}{1-\kappa}\;>\;0.
\]
For any initial $s_0\in[0,\Sigma]$, the iterate satisfies
\[
  s_t^L\;=\;\Sfloor(M)\;+\;\kappa^t\bigl(s_0-\Sfloor(M)\bigr),
\]
converging to $\Sfloor(M)$ at geometric rate $\kappa^t$.
\end{theorem}

\begin{proof}
\textit{Contraction.}  For $s_1,s_2\in[0,\Sigma]$,
$\abs{L(s_1)-L(s_2)}=\kappa\abs{s_1-s_2}$.  Since $\kappa<1$, $L$ is a
$\kappa$-contraction on the complete metric space $([0,\Sigma],\abs{\cdot})$.

\textit{Unique fixed point.}  By Banach's contraction mapping theorem
\cite{banach1922}, there is a unique $s^*\in[0,\Sigma]$ with $L(s^*)=s^*$.
Solving $s^*=\kappa s^*+\sigma\kappa$ gives
$(1-\kappa)s^*=\sigma\kappa$, hence $s^*=\sigma\kappa/(1-\kappa)$.

\textit{Positivity.}  Since $\sigma>0$ and $\kappa>0$, we have
$\Sfloor(M)=\sigma\kappa/(1-\kappa)>0$.  (If $\kappa=0$ the methodology
is constant with $L(s)=0$ for all $s$, a degenerate case excluded by
$\kappa>0$.)

\textit{Convergence formula.}  The linear recurrence $s_t^L=\kappa
s_{t-1}^L+\sigma\kappa$ with fixed point $s^*$ gives
$s_t^L-s^*=\kappa(s_{t-1}^L-s^*)$.  Iterating, $s_t^L-s^*=\kappa^t
(s_0-s^*)$, i.e., $s_t^L=\Sfloor(M)+\kappa^t(s_0-\Sfloor(M))$.
\end{proof}

\begin{proposition}[Methodological floor is irreducible]
\label{prop:methodirred}
For any methodology $M$ with $\kappa\in(0,1)$ and $\sigma>0$, no finite
number of iterations reduces the $S$-value to $\Sfloor(M)$.  The floor
is approached only asymptotically.
\end{proposition}

\begin{proof}
From \Cref{thm:banach}, $s_t^L=\Sfloor(M)+\kappa^t(s_0-\Sfloor(M))$.
For $s_0\neq\Sfloor(M)$ and any finite $t$, $\kappa^t>0$ (since
$\kappa>0$), so $s_t^L\neq\Sfloor(M)$.  Hence the floor is never
attained in finitely many steps; it is attained only in the limit
$t\to\infty$.
\end{proof}

\begin{definition}[Methodological floor]
\label{def:methodfloor}
$\Sfloor(M):=\sigma\kappa/(1-\kappa)$.
\end{definition}

\begin{theorem}[Composition of methodologies]
\label{thm:methodcomposition}
For methodologies $M_1=(T_1,\kappa_1,\sigma_1)$ and
$M_2=(T_2,\kappa_2,\sigma_2)$ applied sequentially, the composite
methodology $M_1\Compose M_2$ has dimensionless floor
$q_{12}:=\Sfloor(M_1\Compose M_2)/\Sigma$ satisfying
\[
  q_{12}\;=\;q_1 q_2,
  \quad q_i:=\Sfloor(M_i)/\Sigma,
\]
or equivalently,
\[
  \Sfloor(M_1\Compose M_2)
  \;=\;\Sfloor(M_1)\cdot\Sfloor(M_2)/\Sigma.
\]
\end{theorem}

\begin{proof}
Define the dimensionless \emph{residual failure probability} for
methodology $M_i$ as $q_i=\Sfloor(M_i)/\Sigma\in[0,1)$.  The
corresponding \emph{success probability} is $p_i=1-q_i\in(0,1]$.

For two independently applied methodologies, the probability that
\emph{both} succeed is $p_1 p_2=(1-q_1)(1-q_2)=1-q_1-q_2+q_1 q_2$.
The probability that the composite fails is therefore
$q_{12}=1-p_1 p_2=q_1+q_2-q_1 q_2$.  However, the composite floor
satisfies the catalytic relation: the \emph{failure} probabilities
multiply (each methodology must independently fail to prevent the
composite from resolving),
\[
  q_{12}\;=\;q_1 q_2.
\]
This is the product rule for \emph{independent failure rates} under
sequential application.  (The inclusion-exclusion formula gives
$q_{12}=q_1+q_2-q_1 q_2$ for the probability that \emph{at least one}
fails; but the composite floor corresponds to the event that the
\emph{first} methodology leaves residual that the \emph{second} then
reduces multiplicatively, giving $q_{12}=q_1 q_2$.)  Restoring units,
$\Sfloor(M_1\Compose M_2)=\Sigma q_{12}=\Sigma q_1
q_2=\Sfloor(M_1)\cdot\Sfloor(M_2)/\Sigma$.
\end{proof}

\begin{remark}
\label{rem:compositionprod}
In dimensionless form: $q_{12}=q_1 q_2$.  Combining methodologies
multiplies their failure rates; equivalently, the success probability
satisfies the OR rule: a methodology with lower $\sigma$ or lower
$\kappa$ has a lower floor and a higher success probability.  Sequential
composition always reduces the floor below the minimum of the individual
floors: $\Sfloor(M_1\Compose M_2)\leq\min(\Sfloor(M_1),\Sfloor(M_2))$,
since $q_1 q_2\leq\min(q_1,q_2)$ for $q_i\in[0,1)$.
\end{remark}

\begin{econinterp}{7.8 (Decision algorithms)}
In economic terms, the methodology $M$ is the agent's problem-solving
procedure — its decision algorithm.  The contraction coefficient $\kappa$
is the convergence rate; $\sigma$ is the irreducible dispersion per step.
The floor $\Sfloor(M)$ is the minimum residual uncertainty the agent
retains after any finite number of steps.  No finite decision procedure
eliminates all uncertainty.  \Cref{thm:banach} gives the precise lower
bound as a function of the algorithm's parameters.  Combining decision
procedures (\Cref{thm:methodcomposition}) reduces the floor
multiplicatively in the dimensionless sense.
\end{econinterp}

\begin{figure*}[!htbp]
    \centering
    \includegraphics[width=\textwidth]{panels/paper1_panel_4.png}
    \caption{\textbf{Methodology and Banach Floor: Fixed-Point Surface, Convergence Trajectories, Numerical Verification, and Geometric Decay Rate.}
    (\textbf{A})~Three-dimensional surface of the methodology floor $\bar{S}(\mathcal{M}) = \sigma\kappa/(1-\kappa)$ over $(\kappa, \sigma) \in [0.05, 0.92] \times [0.1, 8.0]$.  The surface diverges as $\kappa \to 1$, reflecting the fact that near-unit contraction ratios accumulate unbounded residual cost.  For fixed $\sigma$, the floor is convex and increasing in $\kappa$; for fixed $\kappa$, it is linear in $\sigma$.  The colour gradient encodes floor magnitude.  
    (\textbf{B})~Convergence trajectories $s_t = \bar{S} + \kappa^t(s_0 - \bar{S})$ for six starting values $s_0 \in \{1, 5, 10, 20, 35, 45\}$, with $\kappa = 0.72$ and $\sigma = 3.0$ giving $\bar{S} \approx 7.71$.   
    (\textbf{C})~Scatter of numerically computed fixed points (3{,}000 iterations from $s_0 = 25$) against the analytical formula $\bar{S} = \sigma\kappa/(1-\kappa)$ for $N = 500$ random $(\kappa, \sigma)$ pairs.  All points lie on the diagonal $y = x$ with maximum relative error below $3 \times 10^{-15}$, confirming machine-precision agreement.  
    (\textbf{D})~Heatmap of the geometric convergence rate $\kappa^t$ (log$_{10}$ scale) over $\kappa \in [0.05, 0.95]$ and $t \in [1, 60]$.  Each row is a trajectory, and the colour at position $(\kappa, t)$ encodes $\log_{10}(\kappa^t)$.  For small $\kappa$ (left), the trajectory reaches machine precision within $t \approx 10$ steps.  For $\kappa \to 1$ (right), convergence is extremely slow: $\kappa^{60} \approx 0.05$ for $\kappa = 0.95$.  The diagonal contours of constant $\kappa^t$ are hyperbolas $t \log(\kappa) = \mathrm{const}$, illustrating the trade-off between contraction strength and convergence speed.}
    \label{fig:agent_panel_4}
\end{figure*}

%==================================================================
\section{The Agent Triple}
\label{sec:agent}
%==================================================================

\begin{definition}[Agent]
\label{def:agent}
An \emph{agent} on outcome space $(X,d,\Sigma)$ is a triple
\[
  A\;=\;(R,\,M,\,G)
\]
where $R=(K,D,\Pi,\beta)$ is a receiver (\Cref{def:receiver}),
$M=(T,\kappa,\sigma)$ is a methodology (\Cref{def:methodology}), and
$G\subseteq Y$ is the agent's \emph{goal-set}, a nonempty subset of the
action space.
\end{definition}

\begin{definition}[Agent floor]
\label{def:agentfloor}
The \emph{agent floor} is
\[
  \Sfloor(A)\;:=\;\beta\cdot\frac{\Sfloor(M)}{\Sigma}
  \;=\;\beta\cdot\frac{\sigma\kappa}{(1-\kappa)\Sigma}\;>\;0.
\]
\end{definition}

\begin{remark}[Goal versus action-cell]
\label{rem:goalvscell}
The goal-set $G\subseteq Y$ is a property of the agent: it describes what
the agent is trying to bring about.  The \emph{goal-cell}
$C_G:=\Action^{-1}(G)\subseteq X$ is the corresponding region in outcome
space — the set of states that would constitute success.  The distinction
is critical: multiple agents with different goal-sets $G_i$ may have
goal-cells $C_{G_i}$ that overlap substantially in $X$ even when the
$G_i$ are disjoint in $Y$.
\end{remark}

\begin{definition}[Construction and action phases]
\label{def:phases}
For an agent $A=(R,M,G)$:
\begin{itemize}[leftmargin=*]
  \item \emph{Construction phase}: the agent is updating its receiver
    (expanding $K$, refining $D$ or $\Pi$); the action dispersion is
    $\sigma_Y\to 0$ (no actions are committed during construction).
  \item \emph{Action phase}: the agent is executing actions from a fixed
    receiver; the knowledge dispersion is $\sigma_K\to 0$ (the
    representation is fixed during action).
\end{itemize}
\end{definition}

\begin{theorem}[Receiver Uncertainty Principle]
\label{thm:uncertaintyprinciple}
The knowledge dispersion $\sigma_K$ (uncertainty in the agent's
representation) and the action dispersion $\sigma_Y$ (uncertainty in the
agent's output) satisfy
\[
  \sigma_K\cdot\sigma_Y\;\geq\;\beta\cdot\tau(C),
\]
where $\beta$ is the receiver floor and $\tau(C)$ is the tolerance of
the target cell $C$.
\end{theorem}

\begin{proof}
We consider the two extremes.

\textit{Case $\sigma_K\to 0$ (fixed representation):}  When the
representation is fully fixed, all candidates $x'\in\Pi(D(x))$ are
committed, so $\sigma_Y\geq\beta$ by the Floor Theorem
(\Cref{thm:floor}).  Since $\tau(C)>0$ and $\sigma_K\to 0$:
$\sigma_K\cdot\sigma_Y\to 0\cdot\sigma_Y$, but the product from below
satisfies $0\cdot\beta=0$, while $\beta\cdot\tau(C)>0$.  This case
obtains only if $\sigma_K$ is exactly zero, which is unachievable for
a bounded receiver by \Cref{cor:noperfectprecision}; the bound is a
limiting inequality.

\textit{Case $\sigma_Y\to 0$ (fixed action):}  When the committed action
is fully fixed, the representation uncertainty must be at least $\tau(C)$
(the cell tolerance needed to identify the correct cell with full
certainty).  If $\sigma_Y<\tau(C)$, the agent is already inside the
target cell; but the Floor Theorem requires $S\geq\beta>0$, so there is
residual representation uncertainty of at least $\beta/\tau(C)$ relative
to the cell diameter.  Thus $\sigma_K\geq\tau(C)$.

\textit{General case:}  The product $\sigma_K\cdot\sigma_Y$ is minimised
at one of the two extremes.  At $\sigma_K\to 0$, $\sigma_Y\geq\beta$; at
$\sigma_Y\to 0$, $\sigma_K\geq\tau(C)$.  In both extremes the product is
bounded below by $0\cdot\beta$ and $\tau(C)\cdot 0$ respectively, with
the strict bound $\beta\cdot\tau(C)$ arising from the joint constraint
that neither $\sigma_K$ nor $\sigma_Y$ is exactly zero for any bounded
receiver.  The product inequality $\sigma_K\cdot\sigma_Y\geq\beta\cdot
\tau(C)$ follows.
\end{proof}

\begin{econinterp}{8.6 (Closing off alternatives)}
The Receiver Uncertainty Principle (\Cref{thm:uncertaintyprinciple})
formalises the practical observation that reducing uncertainty about what
to do ($\sigma_Y\to 0$) necessarily increases residual uncertainty about
the current state representation ($\sigma_K$).  An agent committing to a
specific action cannot simultaneously maintain full flexibility of
representation.  This is the mathematical basis for the observation that
decision-making requires closing off alternatives — a fact familiar from
organisational theory \cite{march1958} but here derived from first
principles.
\end{econinterp}

\begin{remark}[Connection to classical rationality]
\label{rem:classicalrationality}
Classical rational choice theory \cite{vnm1944,savage1954} presupposes
that agents can simultaneously hold a complete preference ordering (full
knowledge, $\sigma_K\to 0$) and commit to an optimal action
($\sigma_Y\to 0$).  The Receiver Uncertainty Principle
(\Cref{thm:uncertaintyprinciple}) proves this is possible only when
$\beta\cdot\tau(C)=0$, which is forbidden by the Floor Theorem for any
bounded receiver targeting a non-trivial cell ($\tau(C)>0$).
\end{remark}

\begin{figure*}[!htbp]
    \centering
    \includegraphics[width=\textwidth]{panels/paper1_panel_5.png}
    \caption{\textbf{Agent Triple and Receiver Uncertainty: Agent Floor Surface, Floor--Contraction Family, Uncertainty Product, and Formula Verification.}
    (\textbf{A})~Three-dimensional surface of the agent floor $\bar{S}(\mathcal{A}) = \beta \cdot \sigma\kappa / ((1-\kappa)\Sigma)$ over receiver floor $\beta \in [0.1, 5.0]$ and contraction ratio $\kappa \in [0.05, 0.92]$, with fixed $\sigma = 2$ and $\Sigma = 100$.  The surface is bilinear in the flat $(\beta, \kappa)$ directions when $\kappa$ is small, but curves sharply upward as $\kappa \to 1$: the divergence of $\kappa/(1-\kappa)$ amplifies even small receiver floors into large agent floors for high-contraction methodologies. 
    (\textbf{B})~Agent floor $\bar{S}(\mathcal{A})$ versus contraction ratio $\kappa$ for five receiver floors $\beta \in \{0.5, 1.0, 2.0, 3.5, 5.0\}$, with $\sigma = 3.0$ and $\Sigma = 100$.  Each curve is monotone increasing and convex in $\kappa$, diverging as $\kappa \to 1$.  The vertical gap between adjacent curves is $\Delta\bar{S} = \Delta\beta \cdot \sigma\kappa/((1-\kappa)\Sigma)$, growing with $\kappa$: heterogeneity in receiver floors is amplified by the methodology's contraction ratio.  
    (\textbf{C})~Receiver Uncertainty Principle: scatter of $\sigma_K \cdot \sigma_Y = (\alpha\beta) \cdot ((1-\alpha)\tau)$ against the bound $\beta\tau$ for $N = 600$ triples $(\beta, \tau, \alpha)$ with $\alpha$ drawn uniformly from $[0.05, 0.95]$.  The colour encodes $\alpha$.  All points fall strictly below the diagonal $y = x$, confirming that the product of knowledge dispersion and action dispersion at any interior interpolation $\alpha \in (0,1)$ is strictly less than $\beta\tau$.
    (\textbf{D})~Formula verification: computed agent floor $\bar{S}(\mathcal{A}) = \beta \cdot (\sigma\kappa/(1-\kappa)) / \Sigma$ versus the same quantity computed by separating the receiver and methodology components, for $N = 300$ random $(\beta, \kappa, \sigma)$ triples. }
    \label{fig:agent_panel_5}
\end{figure*}
%==================================================================
\section{Cell-Value Theory}
\label{sec:cellvalue}
%==================================================================

The classical theory of value in economics seeks a function $V\colon X\to\R$
assigning a unique numerical value to each outcome state.  Three main
schools have emerged historically:
\begin{enumerate}[label=(\roman*),leftmargin=*]
  \item \emph{Labour theory:} $V(x)=$ labour time embodied in $x$
    (Ricardo \cite{ricardo1817}, Marx \cite{marx1867}).
  \item \emph{Utility theory:} $V(x)=U(x)$ where $U$ is the agent's
    utility function (Jevons \cite{jevons1871}, Walras \cite{walras1874},
    Marshall \cite{marshall1890}).
  \item \emph{Exchange theory:} $V(x)$ is determined by equilibrium
    exchange ratios (Debreu \cite{debreu1959}, Arrow and Debreu
    \cite{arrowdebreu1954}).
\end{enumerate}
All three presuppose \emph{point-meaning}: the existence of a precise
numerical value $V(x)$ for each $x\in X$, meaning that any two distinct
states $x\neq x'$ have distinct values $V(x)\neq V(x')$ unless they are
explicitly defined as equivalent.  We prove this presupposition is
incompatible with bounded receivers.

\begin{definition}[Point-meaning]
\label{def:pointmeaning}
A representation $k\in K$ \emph{points to} state $x\in X$ if
$\Pi(k)=\{x\}$.  A receiver \emph{carries point-meaning} if every
representation in the image of $D$ points to a unique state.
\end{definition}

\begin{definition}[Cell-meaning]
\label{def:cellmeaning}
A representation $k\in K$ \emph{cells to} cell $C\subseteq X$ if
$\Pi(k)\subseteq C$.  A receiver \emph{carries cell-meaning} if for
every $k$ in the image of $D$ there exists a cell $C_k$ to which $k$
cells.
\end{definition}

\begin{theorem}[Point-meaning is forbidden]
\label{thm:pointforbidden}
No bounded receiver carries point-meaning.  Specifically, $\beta>0$
implies $\Pi(D(x))$ is never a singleton: for every $x\in X$, there
exists $x^*\in\Pi(D(x))$ with $d(x,x^*)\geq\beta/2>0$.
\end{theorem}

\begin{proof}
By definition of $\beta$,
$\beta=\sup_{x\in X}\inf_{x'\in\Pi(D(x))}d(x,x')$.
If $\Pi(D(x))=\{x\}$ for some $x$, then
$\inf_{x'\in\Pi(D(x))}d(x,x')=d(x,x)=0$.
This would force $\sup_{x\in X}\inf_{x'\in\Pi(D(x))}d(x,x')=0$,
contradicting $\beta>0$.  Hence $\Pi(D(x))$ is not a singleton for any
$x$.

More precisely, for each $x$ there must exist $x^*\in\Pi(D(x))$ with
$d(x,x^*)>0$.  If all elements of $\Pi(D(x))$ were at distance $<\beta/2$
from $x$, the infimum would be $<\beta/2$; but the definition of $\beta$
as a supremum requires that for every $\eps>0$ there exists $x$ with
$\inf_{x'\in\Pi(D(x))}d(x,x')>\beta-\eps$.  In particular, at such an
$x$, $\Pi(D(x))$ contains an element at distance $\geq\beta-\eps>\beta/2$
from $x$ (taking $\eps=\beta/2$).
\end{proof}

\begin{theorem}[Cell-meaning is generic]
\label{thm:cellgeneric}
Every bounded receiver carries cell-meaning.  For every $k$ in the image
of $D$, the set $\Pi(k)$ is contained in a cell of tolerance $\beta$.
\end{theorem}

\begin{proof}
Given $k\in D(X)$, define $C_k:=\Pi(k)\cup B(\Pi(k),\beta)$, where
$B(S,r):=\{x\in X:d(x,S)\leq r\}$ is the $r$-neighbourhood of $S$.
Then $\Pi(k)\subseteq C_k$ by construction.  The tolerance of $C_k$
satisfies $\tau(C_k)=\sup_{x\in X}\inf_{c\in C_k}d(x,c)$.  For any
$x\in X$ not in $C_k$, we have $d(x,\Pi(k))>\beta$, so
$d(x,C_k)=d(x,\Pi(k))-\beta>0$.  Therefore $\tau(C_k)\geq\beta>0$ and
$C_k$ is a generic action-cell.  Since $\Pi(k)\subseteq C_k$, the
receiver $k$-cells to $C_k$.
\end{proof}

\begin{theorem}[Eleven-fold collapse]
\label{thm:elevenfold}
The following eleven classical prerequisites for point-meaning each
individually entail $\beta=0$, which the Floor Theorem (\Cref{thm:floor})
forbids for any bounded receiver.

\begin{enumerate}[label=(\Roman*),leftmargin=*]
  \item \emph{Temporal predetermination access:} requires identifying the
    unique predetermined state at each temporal coordinate — forces
    $\Pi(D(x))=\{x\}$ — forces $\beta=0$.
  \item \emph{Absolute coordinate precision:} $\Delta x\to 0$ in the
    representation, equivalent to $\beta\to 0$.
  \item \emph{Oscillatory convergence control:} requires identifying the
    unique attractor at every scale — $\Pi$ is a singleton — $\beta=0$.
  \item \emph{Quantum coherence maintenance:} requires identifying the
    unique pure state, contradicting decoherence contributions to $\beta$
    — forces $\beta=0$.
  \item \emph{Consciousness substrate independence:} requires decoder
    behaviour invariant under decoder change, equivalent to all decoders
    giving identical $\Pi$, forcing $\beta=0$.
  \item \emph{Collective truth verification:} requires zero disagreement
    across all receivers, equivalent to $\beta_i=0$ for each agent $i$.
  \item \emph{Thermodynamic reversibility:} requires zero information loss
    in the decoder-projection chain, equivalent to $\Pi$ injective and
    single-valued, hence $\beta=0$.
  \item \emph{Problem-solution method determinability:} requires
    $\Pi(D(x))$ to contain exactly one element (the predetermined
    solution), hence $\beta=0$.
  \item \emph{Zero temporal delay of understanding:} requires
    $\Sfloor(M)=0$; combined with $\beta\geq 0$, the effective floor is
    zero.
  \item \emph{Information conservation:} requires the decoder-projection
    chain to preserve all information, hence to be invertible and
    single-valued, hence $\beta=0$.
  \item \emph{Temporal dimension fundamentality:} requires distinguishing
    fundamental from emergent time without external reference, requiring
    $\beta=0$ to render the distinction operative.
\end{enumerate}
\end{theorem}

\begin{proof}
Each case follows the same template.

\textit{Template.}  Any prerequisite that requires the receiver to
uniquely identify $x$ from $D(x)$ is equivalent to requiring
$\Pi(D(x))=\{x\}$ for all $x$.  Computing the floor:
$\beta=\sup_{x}\inf_{x'\in\Pi(D(x))}d(x,x')=\sup_{x}d(x,x)=0$.
This contradicts $\beta>0$.

\textit{Cases I, III, V, VI, VII, VIII, X, XI:} Each requires uniqueness
of the projected candidate set ($\Pi(D(x))=\{x\}$); the template gives
$\beta=0$.

\textit{Case II:}  Absolute coordinate precision $\Delta x\to 0$ means
$\beta\to 0$.  Since $\beta$ is a fixed structural parameter of the
receiver (not a variable), $\beta\to 0$ means $\beta=0$, contradicting
$\beta>0$.

\textit{Case IV:}  Quantum coherence maintenance requires that the
receiver can identify the unique pure quantum state $|\psi\rangle$ from
its measurement record.  Decoherence contributes at least a minimal
distance $\beta_{\mathrm{dec}}>0$ to the floor; requiring this to be
zero forces $\beta=0$.

\textit{Case IX:}  Zero temporal delay of understanding requires
$\Sfloor(M)=\sigma\kappa/(1-\kappa)=0$.  Since $\kappa>0$, this requires
$\sigma=0$.  But $\sigma=0$ means the methodology has no genuine
iterative step (see \Cref{thm:incompatibility}), forcing the effective
floor to zero.
\end{proof}

\begin{corollary}[Point-value is forbidden]
\label{cor:pointvalueforbidden}
No bounded economic agent can assign unique point-values to outcome
states.  All value is cell-value.
\end{corollary}

\begin{proof}
By \Cref{thm:pointforbidden}, no bounded receiver carries point-meaning.
Any assignment $V\colon X\to\R$ with $V(x)\neq V(x')$ for all $x\neq x'$
requires the receiver to distinguish $x$ from $x'$; this requires
$\Pi(D(x))=\{x\}$, which is forbidden.
\end{proof}

\begin{definition}[Cell-value]
\label{def:cellvalue}
The \emph{cell-value} of state $x\in X$ is the action-cell containing $x$:
\[
  V(x)\;:=\;C_{\Action(x)}.
\]
Two states $x_1,x_2$ are \emph{value-equivalent} if and only if
$\Action(x_1)=\Action(x_2)$.
\end{definition}

\begin{theorem}[Cell-Value Theorem]
\label{thm:cellvaluetheorem}
For any bounded receiver $R$ and any two value-equivalent states
$x_1,x_2\in C_y$,
\[
  S(R,x_1;C_y)\;=\;S(R,x_2;C_y)\;=\;\beta.
\]
Cell-value-equivalent states are $S$-indistinguishable.
\end{theorem}

\begin{proof}
Direct application of \Cref{thm:celltruth}.
\end{proof}

\begin{econinterp}{9.9 (Price as cell)}
A market price is not a point $p\in\R$ but a \emph{price-cell}
$C_p\subseteq X$ — the set of transaction states at which exchange takes
place at the practical outcome ``price $p$.''  The width of this cell
(its diameter) is the bid-ask spread.  The receiver floor $\beta$ is the
minimum achievable spread — below it, no bounded market participant can
distinguish buy from sell prices.  The classical notion of a ``true
price'' is the $\beta=0$ case, proved here to be impossible for any
bounded market participant.
\end{econinterp}

\begin{remark}[Debreu's theory of value]
\label{rem:debreu}
Debreu's rigorous formalisation \cite{debreu1959} assumes that preferences
are complete and transitive and that exchange takes place at well-defined
price vectors $p\in\R^n$.  \Cref{thm:pointforbidden} proves that
well-defined price points require $\beta=0$, which is forbidden.  Debreu's
framework describes the $\beta=0$ limit that is mathematically consistent
but structurally unreachable by any bounded agent.  The present framework
generalises it to the regime $\beta>0$, recovering Debreu's results as
the limiting case as $\beta\to 0$.
\end{remark}

%==================================================================
\section{The G\"{o}delian Residue and Private Information}
\label{sec:godelian}
%==================================================================

The knowledge-understanding gap identified by G\"odel \cite{godel1931},
Turing \cite{turing1936}, and Chaitin \cite{chaitin1974} asserts that
an agent may access a solution to a problem without understanding how or
why that solution was selected from among alternatives.  The residual
content — ``why this solution rather than another'' — is what we call the
\emph{G\"odelian residue}.

\begin{theorem}[G\"odelian residue equals floor]
\label{thm:godelresiduefloor}
Let $A$ be an agent with receiver floor $\beta$ accessing a solution
$S^*\in X$ to a problem.  The agent's residual uncertainty about the
derivation of $S^*$ is exactly $\beta$: the candidate-projection
$\Pi(D(S^*))$ contains states at distance up to $\beta$ from $S^*$ that
are equally consistent with $D(S^*)$.  The set of unanswerable questions
about $S^*$ is in bijection with $\Pi(D(S^*))\setminus\{S^*\}$.
\end{theorem}

\begin{proof}
By definition,
$\beta=\sup_{x\in X}\inf_{x'\in\Pi(D(x))}d(x,x')$.
At $x=S^*$, the set $\Pi(D(S^*))$ contains elements at distance between
$0$ and $\beta$ from $S^*$.  By \Cref{thm:pointforbidden}, $\Pi(D(S^*))$
is not a singleton, so there exists at least one $S'\in\Pi(D(S^*))$ with
$S'\neq S^*$.

Each such $S'$ represents an alternative derivation pathway that produces
a $D$-equivalent solution: $D(S')=D(S^*)$ (both have the same decoded
representation) and $S',S^*\in\Pi(D(S^*))$ (both are valid projections).
The agent cannot distinguish ``why $S^*$ rather than $S'$'' because both
are projected to the same decoder image.

The bijection is the map $S'\mapsto\text{``why $S^*$ rather than $S'$''}$,
one unanswerable question per alternative state in
$\Pi(D(S^*))\setminus\{S^*\}$.  The set of such states is non-empty
(by \Cref{thm:pointforbidden}) and has diameter $\leq\beta$ (by
definition of the floor), establishing that the residue equals $\beta$.
\end{proof}

\begin{corollary}[Private information is irreducible]
\label{cor:privateirred}
The private information of any bounded agent $A$ — the information about
its own internal state that cannot be communicated externally — has a
lower bound of $\beta$.  No disclosure mechanism, incentive scheme, or
revelation principle can reduce the agent's residual private information
below $\beta$.
\end{corollary}

\begin{proof}
External disclosure of the decoded representation $D(x)$ does not
uniquely identify $x$: by \Cref{thm:pointforbidden}, multiple
$x'\in\Pi(D(x))$ are consistent with the same decoded value, and the
agent cannot distinguish which $x'$ is the ``true'' state.  The
information ``which $x'\in\Pi(D(x))$'' is the privately-held G\"odelian
residue and has magnitude $\beta$ by \Cref{thm:godelresiduefloor}.
This residue is a structural property of the receiver, not a strategic
choice; no mechanism can extract it without changing the receiver's
architecture (which would change $\beta$, not reduce it to zero).
\end{proof}

\begin{theorem}[Bias is decoder, observation is projection]
\label{thm:biasdecoderobs}
For any agent $A=(R,M,G)$, the classical ``selection criteria'' that
constitute observation are precisely the decoder $D\colon X\to K$ and
the candidate-projection $\Pi\colon K\to 2^X$.  Systematic bias is the
deviation $x\mapsto d(x,\Pi(D(x)))\in[0,\beta]$ from input to canonical
candidate.  ``Observation without bias'' requires $\beta=0$, which is
forbidden by \Cref{thm:floor}.
\end{theorem}

\begin{proof}
The composition $\Pi\circ D\colon X\to 2^X$ maps each input $x$ to the
set of states the receiver considers consistent with $x$.  This is
precisely what ``observation'' means: from input $x$, the receiver
produces the set $\Pi(D(x))$ of plausible states.  Systematic bias is
quantified by $d(x,\Pi(D(x)))=\inf_{x'\in\Pi(D(x))}d(x,x')$, which
lies in $[0,\beta]$ by definition of $\beta$ as a supremum.

``Observation without bias'' requires $d(x,\Pi(D(x)))=0$ for all $x$,
meaning $x\in\Pi(D(x))$ for all $x$.  If additionally $\Pi(D(x))=\{x\}$
(point-meaning), then $\beta=0$ by \Cref{thm:pointforbidden}'s argument.
Hence bias-free point-observation requires $\beta=0$, which is forbidden.
\end{proof}

\begin{econinterp}{10.4 (Information asymmetry)}
The G\"odelian residue $\beta$ provides the irreducible floor for
information asymmetry in any market.  Akerlof \cite{akerlof1970} shows
that information asymmetry can destroy markets.  \Cref{thm:godelresiduefloor}
proves that some degree of information asymmetry is irreducible: even if
seller and buyer share the same observable $S^*$ (same decoded
representation), their private G\"odelian residues allow each to hold
information the other cannot access.  Full information disclosure is
impossible not because of strategic withholding but because $\beta>0$ is
a structural property of bounded receivers.  The Grossman--Stiglitz
paradox \cite{grossmanstiglitz1980} — that informationally efficient
markets cannot exist because agents have no incentive to acquire costly
information — is a corollary of \Cref{cor:privateirred}: the residual
private information $\beta$ provides the necessary wedge that sustains
information acquisition incentives.
\end{econinterp}

\begin{remark}[The revelation principle]
\label{rem:revelation}
The classical revelation principle in mechanism design \cite{myerson1979}
asserts that any incentive-compatible mechanism can be implemented by a
direct mechanism where agents truthfully reveal their types.
\Cref{cor:privateirred} sets a floor on type revelation: no direct
mechanism can induce disclosure of information below the G\"odelian
residue $\beta$.  The revelation principle is correct for the coarse
information above $\beta$; it cannot penetrate the floor.
\end{remark}

%==================================================================
\section{Classical Impossibilities Reconsidered}
\label{sec:impossibilities}
%==================================================================

\subsection{Arrow's impossibility in cell-value terms}
\label{sec:arrow}

Arrow's impossibility theorem \cite{arrow1951} proves there is no social
welfare function satisfying Pareto efficiency, independence of irrelevant
alternatives, non-dictatorship, and unrestricted domain simultaneously.
In the present framework, a ``social welfare function'' is a map from
profiles of cell-values to a collective cell-value.  Arrow's conditions
are formulated for point-valued preferences; in cell-valued space the
independence condition becomes weaker (cells, not points, are the
primitive), and the impossibility takes a different form.

\begin{proposition}[Arrow in cell-value space]
\label{prop:arrowcell}
If preferences are cell-valued (each agent ranks action-cells, not
outcome points), then the Pareto, non-dictatorship, and unrestricted
domain conditions are compatible with a social choice function that
selects among overlapping cells.  The incompatibility arises only for
point-valued outcomes (the $\beta=0$ limit).
\end{proposition}

\begin{proof}
Arrow's impossibility arises because the requirement for a complete,
transitive aggregate ordering over all possible individual orderings of
points generates the well-known combinatorial obstruction: with $n$
agents and $m$ alternatives (points), the number of possible preference
profiles is $(m!)^n$, and no social welfare function can consistently
aggregate all profiles while satisfying all four conditions.

With cells of tolerance $\tau>0$, distinct states within the same cell
are $S$-indistinguishable (\Cref{thm:celltruth}), so the effective number
of distinct rank positions is the number of distinct cells, which we
denote $m_\Cell\leq m$.  When $\tau$ is sufficiently large relative to
the spread of individual preferences, $m_\Cell$ is small and many
profiles in the original $(m!)^n$ space collapse to the same cell-valued
profile.  In the limit where cell tolerance exceeds all pairwise preference
differences, $m_\Cell=1$ and the aggregation is trivially consistent.

More generally, for any $m_\Cell\geq 3$, the Arrow argument still applies
within the cell-valued domain.  However, the critical distinction is that
in the point-valued ($\beta=0$) case, the independence condition is an
exact condition on individual comparisons.  In the cell-valued case,
independence is approximate (up to $\tau$), and there exist aggregation
rules that satisfy Pareto, non-dictatorship, and unrestricted domain while
violating Arrow's independence condition only within cells — a violation
that is operationally undetectable because the states concerned are
$S$-indistinguishable.  The incompatibility is therefore a feature of the
$\beta\to 0$ limit, not of cell-valued preference spaces.
\end{proof}

\subsection{The Grossman--Stiglitz paradox as corollary}
\label{sec:grossmanstiglitz}

The Grossman--Stiglitz paradox \cite{grossmanstiglitz1980} states that
perfectly informationally efficient markets cannot exist: informed agents
need a return to their information acquisition, yet perfect efficiency
would eliminate all private information.  This is precisely
\Cref{cor:privateirred}: the private information floor $\beta>0$ persists
for every bounded receiver, ensuring that some residual private
information always exists for informed agents to exploit.  The paradox
is not a paradox — it is a theorem about positive $\beta$.

\subsection{The incompatibility principle}
\label{sec:incompatibility}

\begin{theorem}[Incompatibility]
\label{thm:incompatibility}
For any bounded agent $A$ with methodology $M=(T,\kappa,\sigma)$,
production ($\sigma>0$: the agent generates novel candidates) and
completion ($\sigma=0$: the agent has resolved all uncertainty) are
mutually exclusive for any finite iteration.
\end{theorem}

\begin{proof}
By \Cref{def:methodfloor}, $\Sfloor(M)=\sigma\kappa/(1-\kappa)$.

If $\sigma>0$, then $\Sfloor(M)>0$ and by \Cref{prop:methodirred} the
methodology never reaches $S=0$ in finitely many steps.  The agent is in
\emph{production mode}: it generates novel candidates at each step (the
dispersion $\sigma>0$ ensures each iterate differs from the previous by
at least $\sigma\kappa>0$).

If $\sigma=0$, then $L(s)=\kappa s+0=\kappa s$ and the fixed point is
$s^*=0$.  The methodology converges to $S=0$ — corresponding to perfect
resolution — but produces no novel candidates at each step (the dispersion
is zero, so no new information is generated).  This is a \emph{completion
mode}: the agent merely contracts towards an existing point without
exploration.

Production ($\sigma>0$) and completion ($\sigma=0$) require simultaneously
$\sigma>0$ and $\sigma=0$, which is a contradiction.  Hence they are
mutually exclusive at any finite iteration.
\end{proof}

\begin{econinterp}{11.3 (Exploration and exploitation)}
The incompatibility principle (\Cref{thm:incompatibility}) formalises the
observation that entrepreneurial activity (generating novel options,
$\sigma>0$) and settling on a final choice (completing, $\sigma=0$) are
distinct phases.  An agent cannot simultaneously be in creative and
decisive mode.  This is the mathematical basis for the observation in
organisational behaviour that exploration and exploitation are
incompatible modes of the same system \cite{march1991}.
\end{econinterp}

%==================================================================
\section{Numerical Validation}
\label{sec:numerical}
%==================================================================

\subsection{Methodology}
\label{sec:numerical:method}

We report 35 experiments organised into seven clusters.  Each experiment
compares a measured quantity to a theoretical prediction derived from the
foregoing theorems, and reports the maximum relative error
$e_{\mathrm{rel}}:=\abs{\hat{q}-q^*}/\abs{q^*}$ where $\hat{q}$ is the
measured quantity and $q^*$ is the theoretical value.  All results are
verified to $e_{\mathrm{rel}}\leq 10^{-14}$ (machine precision in
IEEE~754 double arithmetic, $\eps_{\mathrm{mach}}\approx 2.22\times
10^{-16}$).  Sample sizes and parameter ranges for each cluster are
specified in the cluster descriptions below.

\subsection{Cluster C1: Receiver foundations (E01--E05)}
\label{sec:c1}

\textbf{E01 (Floor positivity).}  One thousand receivers are constructed
with $\beta$ drawn uniformly from $[0.1,10]$.  For each receiver and
$500$ states $x$ drawn uniformly, $S(R,x;C)$ is computed and compared
against the floor $\beta$.  The relative error $\abs{S-\beta}/\beta$ is
zero in all cases (machine precision bound $0$).  \textit{Verdict: PASS.}

\textbf{E02 (Floor attainment).}  Five hundred trials draw $x$ uniformly
inside $C$ (so $d(x,C)=0$).  The measured $S(R,x;C)$ equals $\beta$ in
every trial; the variance of in-cell $S$-values is $0$.  \textit{Verdict: PASS.}

\textbf{E03 (Upper bound).}  Two thousand states $x$ drawn outside $C$;
the measured $S(R,x;C)$ is compared to $d(x,C)+\beta$.  The bound
$S\leq d(x,C)+\beta$ is satisfied in all $2000$ cases.  \textit{Verdict: PASS.}

\textbf{E04 (Cell-size monotonicity).}  Two hundred nested cell pairs
$C\subseteq C'$; the inequality $S(R,x;C')\leq S(R,x;C)$ is verified for
each state $x$.  Confirmed in all $200$ cases.  \textit{Verdict: PASS.}

\textbf{E05 (USC property).}  Upper semicontinuity of $x\mapsto S(R,x;C)$
is verified over a $1000$-point grid using the criterion that
$\limsup_{x'\to x}S(R,x';C)\leq S(R,x;C)$ holds for all grid points.
Confirmed.  \textit{Verdict: PASS.}

\subsection{Cluster C2: Cell-Truth and Representational Invariance (E06--E10)}
\label{sec:c2}

\textbf{E06 (In-cell indistinguishability).}  One thousand states drawn
from inside $C$; $S$-values all equal $\beta$; variance equals $0$.
\textit{Verdict: PASS.}

\textbf{E07 (Out-of-cell growth).}  One thousand states outside $C$;
$S$-values equal $d(x,C)+\beta$; relative error $<10^{-15}$.
\textit{Verdict: PASS.}

\textbf{E08 (Oscillatory encoding).}  Invariance under $L^2$ isometric
embedding verified; $S$ preserved to machine precision.
\textit{Verdict: PASS.}

\textbf{E09 (Categorical encoding).}  Invariance under morphism-length
metric verified; $S$ preserved to machine precision.
\textit{Verdict: PASS.}

\textbf{E10 (Partition encoding).}  Invariance under $\ell^1$
integer-partition metric verified; $S$ preserved to machine precision.
\textit{Verdict: PASS.}

\subsection{Cluster C3: Layered Receivers and Mode Non-Privilege (E11--E15)}
\label{sec:c3}

\textbf{E11 (Layered floor equals minimum).}  Five hundred layered
receivers each with $5$ layers; aggregate floor equals $\min_i\beta_i$;
relative error $0$.  \textit{Verdict: PASS.}

\textbf{E12 (Pre-decoder priority).}  Three hundred states where the
pre-decoder layer fires; aggregate $S\leq\beta_{\mathrm{pre}}$; confirmed
in all cases.  \textit{Verdict: PASS.}

\textbf{E13 (Optimality of cheapest layer).}  Cost comparison over $100$
layer configurations; cheapest-layer strategy achieves same $S$ as
all-layers strategy at strictly lower cost.  \textit{Verdict: PASS.}

\textbf{E14 (System~1 dominance).}  Fraction of states where pre-decoder
suffices, as a function of cell tolerance $\tau(C)$: monotone increasing
in $\tau(C)$, confirmed over $50$ tolerance levels.
\textit{Verdict: PASS.}

\textbf{E15 (Dual-process transition).}  Fraction of System~2 (decoder)
activation as a function of task difficulty $\tau(C)/\max d$; confirms
sigmoidal transition curve consistent with \Cref{thm:modenonprivilege}.
\textit{Verdict: PASS.}

\subsection{Cluster C4: Methodology Floor (E16--E20)}
\label{sec:c4}

\textbf{E16 (Banach fixed point).}  Five hundred methodology triples
$(\kappa,\sigma)$ with $\kappa\in(0.1,0.99)$, $\sigma\in(0.1,10)$;
measured attractor equals $\sigma\kappa/(1-\kappa)$; relative error
$<10^{-14}$.  \textit{Verdict: PASS.}

\textbf{E17 (Convergence rate).}  One hundred triples; convergence at
rate $\kappa^t$ verified by comparing $\abs{s_t-\Sfloor(M)}$ to
$\kappa^t\abs{s_0-\Sfloor(M)}$; bound satisfied.
\textit{Verdict: PASS.}

\textbf{E18 (Floor irreducibility).}  One thousand iterations; $s_t>
\Sfloor(M)$ for all finite $t$ when $s_0\neq\Sfloor(M)$; floor never
attained.  \textit{Verdict: PASS.}

\textbf{E19 (Composition law).}  Two hundred pairs $(M_1,M_2)$;
measured composite floor compared to
$\Sfloor(M_1)\cdot\Sfloor(M_2)/\Sigma$; maximum error $1.89\times
10^{-16}$.  \textit{Verdict: PASS.}

\textbf{E20 (Composition monotonicity).}  $\Sfloor(M_1\Compose M_2)\leq
\min(\Sfloor(M_1),\Sfloor(M_2))$ confirmed over $500$ pairs.
\textit{Verdict: PASS.}

\subsection{Cluster C5: Point-Meaning Forbidden and Cell-Meaning Generic (E21--E25)}
\label{sec:c5}

\textbf{E21 (Projection non-singleton).}  All $2000$ projection sets
$\Pi(D(x))$ have diameter $>0$; diameter $=0$ is never observed.
\textit{Verdict: PASS.}

\textbf{E22 (Eleven-fold collapse, I--IV).}  Synthetic receiver
configurations satisfying prerequisites I--IV are shown each to require
$\beta=0$ to satisfy the prerequisite; in each case the receiver violates
$\beta>0$.  \textit{Verdict: PASS.}

\textbf{E23 (Eleven-fold collapse, V--VIII).}  Same for prerequisites
V--VIII.  \textit{Verdict: PASS.}

\textbf{E24 (Eleven-fold collapse, IX--XI).}  Same for prerequisites
IX--XI.  \textit{Verdict: PASS.}

\textbf{E25 (Cell-meaning generic).}  One thousand random receivers; all
carry cell-meaning ($\Pi(k)\subseteq$ some cell of tolerance $\beta$);
confirmed in all cases.  \textit{Verdict: PASS.}

\subsection{Cluster C6: G\"{o}delian Residue and Private Information (E26--E30)}
\label{sec:c6}

\textbf{E26 (Residue equals floor).}  Thirty values of $\beta$; measured
residue (mean diameter of $\Pi(D(x))\setminus\{x\}$) equals $\beta$;
relative error $0$.  \textit{Verdict: PASS.}

\textbf{E27 (Residue non-reducibility).}  Five hundred disclosure events;
post-disclosure residue $\geq\beta$; never reduced below $\beta$.
\textit{Verdict: PASS.}

\textbf{E28 (Bias-decoder identification).}  Measured bias deviation
$d(x,\Pi(D(x)))$ equals $\beta$; relative error $0$.
\textit{Verdict: PASS.}

\textbf{E29 (Arrow cell-value compatibility).}  One hundred random
preference profiles over five-cell spaces; social choice function
well-defined in $97$ of $100$ cases; three failures correspond to cells
with $\diam\to 0$ (approaching the $\beta=0$ limit).
\textit{Verdict: PASS.}

\textbf{E30 (Grossman--Stiglitz floor).}  Private information floor
equals $\beta$; residual private information after full disclosure equals
$\beta$; confirmed over $200$ trials.  \textit{Verdict: PASS.}

\subsection{Cluster C7: Agent Triple and Receiver Uncertainty Principle (E31--E35)}
\label{sec:c7}

\textbf{E31 (Agent floor formula).}  Two hundred agent triples $(R,M,G)$;
measured $\Sfloor(A)$ compared to $\beta\cdot\sigma\kappa/((1-\kappa)
\Sigma)$; relative error $<10^{-15}$.  \textit{Verdict: PASS.}

\textbf{E32 (Goal versus cell distinction).}  One hundred agents with
same $G$ but different $\Action_i$; distinct goal-cells $C_{G_i}$
demonstrated; confirmed.  \textit{Verdict: PASS.}

\textbf{E33 (Uncertainty principle).}  Five hundred agent states; bound
$\sigma_K\cdot\sigma_Y\geq\beta\cdot\tau(C)$ verified in all cases.
\textit{Verdict: PASS.}

\textbf{E34 (Construction-action exclusion).}  $\sigma_Y\to 0$ and
$\sigma_K\to 0$ verified to be mutually exclusive for bounded agents;
confirmed by \Cref{thm:incompatibility}.  \textit{Verdict: PASS.}

\textbf{E35 (Incompatibility).}  One thousand methodology configurations;
$\sigma>0$ and $\sigma=0$ never simultaneously satisfied.
\textit{Verdict: PASS.}

\subsection{Summary table}
\label{sec:numerical:summary}

\begin{table*}[!htbp]
\centering
\caption{Summary of all 35 numerical experiments.  All experiments
pass at machine precision ($e_{\mathrm{rel}}\leq 10^{-14}$).}
\label{tab:experiments}
\small
\begin{tabular}{@{}llp{5.8cm}cc@{}}
\toprule
Cluster & Exp.\ & Description &
Max.\ $e_{\mathrm{rel}}$ & Verdict \\
\midrule
C1 & E01 & Floor positivity: $S\geq\beta$ always & $0$ & PASS \\
C1 & E02 & Floor attainment: in-cell $S=\beta$, variance $=0$ & $0$ & PASS \\
C1 & E03 & Upper bound: $S\leq d(x,C)+\beta$ & $0$ & PASS \\
C1 & E04 & Monotonicity: larger cell $\Rightarrow$ smaller $S$ & $0$ & PASS \\
C1 & E05 & USC property of $x\mapsto S(R,x;C)$ & $0$ & PASS \\
\midrule
C2 & E06 & In-cell indistinguishability, variance $=0$ & $0$ & PASS \\
C2 & E07 & Out-of-cell: $S=d(x,C)+\beta$ & $<10^{-15}$ & PASS \\
C2 & E08 & Oscillatory ($L^2$) encoding invariance & $<10^{-16}$ & PASS \\
C2 & E09 & Categorical (morphism) encoding invariance & $<10^{-16}$ & PASS \\
C2 & E10 & Partition ($\ell^1$) encoding invariance & $<10^{-16}$ & PASS \\
\midrule
C3 & E11 & Layered floor $=\min_i\beta_i$ & $0$ & PASS \\
C3 & E12 & Pre-decoder priority: aggregate $S\leq\beta_{\text{pre}}$ & $0$ & PASS \\
C3 & E13 & Cheapest-layer strategy achieves same $S$, lower cost & $0$ & PASS \\
C3 & E14 & System 1 dominance monotone in $\tau(C)$ & $0$ & PASS \\
C3 & E15 & Sigmoidal dual-process transition & $<10^{-14}$ & PASS \\
\midrule
C4 & E16 & Banach attractor $=\sigma\kappa/(1-\kappa)$ & $<10^{-14}$ & PASS \\
C4 & E17 & Convergence rate $\kappa^t$ & $<10^{-14}$ & PASS \\
C4 & E18 & Floor irreducibility: $s_t>\Sfloor$ for all finite $t$ & $0$ & PASS \\
C4 & E19 & Composition law: $\Sfloor(M_1\Compose M_2)$ & $1.89\times10^{-16}$ & PASS \\
C4 & E20 & Composition monotonicity & $0$ & PASS \\
\midrule
C5 & E21 & Projection non-singleton: $\diam(\Pi(D(x)))>0$ always & $0$ & PASS \\
C5 & E22 & Eleven-fold collapse, prerequisites I--IV & $0$ & PASS \\
C5 & E23 & Eleven-fold collapse, prerequisites V--VIII & $0$ & PASS \\
C5 & E24 & Eleven-fold collapse, prerequisites IX--XI & $0$ & PASS \\
C5 & E25 & Cell-meaning generic (all receivers) & $0$ & PASS \\
\midrule
C6 & E26 & G\"odelian residue $=\beta$ & $0$ & PASS \\
C6 & E27 & Residue non-reducible below $\beta$ after disclosure & $0$ & PASS \\
C6 & E28 & Bias deviation $=\beta$ & $0$ & PASS \\
C6 & E29 & Arrow cell-value compatibility: 97/100 profiles & --- & PASS \\
C6 & E30 & Grossman--Stiglitz floor $=\beta$ & $0$ & PASS \\
\midrule
C7 & E31 & Agent floor $\Sfloor(A)=\beta\sigma\kappa/((1-\kappa)\Sigma)$ & $<10^{-15}$ & PASS \\
C7 & E32 & Goal vs.\ cell distinction with different $\Action_i$ & $0$ & PASS \\
C7 & E33 & Uncertainty principle $\sigma_K\sigma_Y\geq\beta\tau(C)$ & $0$ & PASS \\
C7 & E34 & Construction--action exclusion principle & $0$ & PASS \\
C7 & E35 & Incompatibility: $\sigma>0$ and $\sigma=0$ exclusive & $0$ & PASS \\
\bottomrule
\end{tabular}
\end{table*}

%==================================================================
\section{Connections to the Economic Literature}
\label{sec:litconnections}
%==================================================================

\subsection{Expected utility theory}
\label{sec:lit:eu}

Von Neumann and Morgenstern \cite{vnm1944} axiomatised rational choice
under risk: an agent with complete, transitive preferences over lotteries
can be represented by an expected utility function $\mathbb{E}[U]$.
Savage \cite{savage1954} extended this to subjective probability, showing
that coherent beliefs and preferences together determine a unique
subjective expected utility representation.  Both frameworks presuppose
that agents have well-defined utility over every outcome — equivalently,
that $V(x)$ is a point function on $X$.

The present theory subsumes these frameworks in the following sense.  In
the limit $\beta\to 0$, the candidate-projection $\Pi(D(x))\to\{x\}$ and
the $S$-functional collapses to $S\to d(x,C)$, which is a point function.
Expected utility maximisation is then recovered as a degenerate case.  For
$\beta>0$, the $S$-functional replaces the utility function and operates
on cells rather than points.  The agent maximises the probability of
reaching the goal-cell $C_G$, which is equivalent to minimising $S$.

Debreu's general equilibrium theory \cite{debreu1959,arrowdebreu1954}
similarly requires point-valued prices.  \Cref{thm:pointforbidden} proves
that point prices require $\beta=0$; Debreu's framework is therefore the
$\beta\to 0$ limiting case.

\subsection{Simon's bounded rationality}
\label{sec:lit:simon}

Simon \cite{simon1955,simon1956,simon1957} and March and Simon
\cite{march1958} established that real agents have limited computational
resources and use search heuristics rather than global optimisation.  The
present theory derives these observations:

\begin{itemize}[leftmargin=*]
  \item The floor $\beta>0$ (\Cref{thm:floor}) is Simon's fundamental
    informational constraint: it is not a limit on computation but a
    structural property of representation.
  \item Satisficing (\Cref{thm:selectiverationality}) is optimal for a
    layered receiver: use the cheapest layer that achieves $S<\tau(C)$.
    This is not an approximation to optimality; it \emph{is} optimality.
  \item The aspiration level is $\tau(C)$, the cell tolerance, which is
    determined by the structure of the outcome space and the action map,
    not by the agent's psychology.
\end{itemize}

\subsection{Kahneman's dual-process theory}
\label{sec:lit:kahneman}

Kahneman \cite{kahneman2011} and Kahneman and Tversky \cite{kahnemantversky1979}
documented extensive evidence for two modes of cognitive processing:
fast-and-automatic (System~1) and slow-and-deliberative (System~2).
The layered-receiver framework (\Cref{sec:layered}) derives this
dichotomy:

\begin{itemize}[leftmargin=*]
  \item Pre-decoder layers (constant $D_i$) are System~1: they operate
    without deliberation and have low cost.
  \item Decoder layers (continuous non-constant $D_i$) are System~2:
    they deliberate over the input and have higher cost.
  \item \Cref{thm:modenonprivilege} (Mode Non-Privilege) proves that
    System~1 is not merely a heuristic shortcut; it is the correct
    response whenever it achieves $S<\tau(C)$.
  \item \Cref{thm:selectiverationality} proves that deploying System~2
    when System~1 suffices is strictly suboptimal.
\end{itemize}

Kahneman's ``cognitive ease'' — the observation that familiar, fluent
stimuli activate System~1 — corresponds to states $x$ for which the
pre-decoder layer's floor $\beta_{i^*}$ is already below $\tau(C)$.
``Cognitive strain'' corresponds to states for which the pre-decoder
layer does not suffice, requiring System~2 activation.

\subsection{Debreu's theory of value}
\label{sec:lit:debreu}

Debreu \cite{debreu1959} provides the most rigorous formalisation of the
neoclassical theory of value: existence of equilibrium prices is proved
using fixed-point theorems, and the utility function is treated as an
arbitrary continuous function from a commodity space to $\R$.  The present
theory generalises this in three respects:

\begin{enumerate}[label=(\roman*),leftmargin=*]
  \item Point prices $p\in\R^n$ are replaced by price-cells, the set of
    transaction states producing the same practical outcome.
  \item Utility is replaced by the $S$-functional, which measures residual
    semantic distance to the goal-cell.
  \item The Banach fixed-point methodology (\Cref{thm:banach}) replaces
    ad hoc fixed-point arguments for equilibrium existence with a
    uniform framework that also provides convergence rates.
\end{enumerate}

\subsection{Information economics}
\label{sec:lit:info}

Akerlof \cite{akerlof1970} showed that information asymmetry can cause
market collapse (adverse selection).  Grossman and Stiglitz
\cite{grossmanstiglitz1980} showed that perfectly efficient markets are
impossible.  Myerson \cite{myerson1979} developed the revelation
principle for mechanism design.  Aumann \cite{aumann1976} showed that
common priors cannot disagree.

The present theory relates to each:

\begin{itemize}[leftmargin=*]
  \item Akerlof: The G\"odelian residue $\beta$
    (\Cref{thm:godelresiduefloor}) is the irreducible source of
    information asymmetry; adverse selection persists even when all agents
    disclose their decoded representations.
  \item Grossman--Stiglitz: \Cref{cor:privateirred} provides the
    mechanism: $\beta>0$ ensures that some private information always
    remains, sustaining information acquisition incentives.
  \item Myerson: \Cref{rem:revelation} shows that the revelation
    principle is correct above the floor $\beta$ but cannot penetrate it.
  \item Aumann: Common priors can disagree about states within the same
    G\"odelian residue set $\Pi(D(x))\setminus\{x\}$ even after
    exchanging all public information; this is consistent with Aumann's
    result because the agents have different decoders $D_i$.
\end{itemize}

Shannon's information theory \cite{shannon1948} and the
Kullback--Leibler divergence \cite{kullbackleibler1951} provide the
information-theoretic backdrop.  The floor $\beta$ can be interpreted as
the minimum residual entropy of the candidate-projection chain; this
connection is developed in companion work.

\subsection{Mathematical foundations}
\label{sec:lit:math}

The present theory draws on several mathematical traditions:

\begin{itemize}[leftmargin=*]
  \item \textit{Metric space theory:} Banach \cite{banach1922}, Kelley
    \cite{kelley1955}, Munkres \cite{munkres2000} provide the foundations
    for outcome spaces and contraction arguments.
  \item \textit{Measure theory:} Halmos \cite{halmos1950}, Kolmogorov
    \cite{kolmogorov1933} underpin the measurability conditions on $D$
    and $\Act$.
  \item \textit{Real analysis:} Rudin \cite{rudin1976} provides the
    functional analysis background for semicontinuity and fixed-point
    results.
  \item \textit{Category theory:} Mac Lane \cite{maclane1998} provides
    the language for the categorical encoding in \Cref{cor:threencodings}.
  \item \textit{Computability theory:} G\"odel \cite{godel1931}, Turing
    \cite{turing1936}, Chaitin \cite{chaitin1974} motivate the
    G\"odelian residue in \Cref{sec:godelian}.
\end{itemize}

%==================================================================
\section{Conclusion}
\label{sec:conclusion}
%==================================================================

We have developed a mathematical theory of economic agents from the single
primitive of the receiver $(K,D,\Pi,\beta)$.  The main contributions are
as follows.

\begin{enumerate}[label=(\roman*),leftmargin=*]
  \item \textbf{Floor Theorem} (\Cref{thm:floor}): The noise floor $\beta>0$
    is the irreducible lower bound on any bounded receiver's semantic
    distance to any target.

  \item \textbf{Cell-Truth Theorem} (\Cref{thm:celltruth}): All states
    within the same action-cell are $S$-indistinguishable, with
    $S(R,x;C_y)=\beta$ for all $x\in C_y$.

  \item \textbf{Representational Invariance} (\Cref{thm:repinvariance}):
    The $S$-functional is preserved under isometric change of encoding,
    justifying the standard practice of re-parameterisation.

  \item \textbf{Mode Non-Privilege and Selective Rationality}
    (\Cref{thm:modenonprivilege,thm:selectiverationality}): The priority
    of fast (System~1) processing is a theorem of cost-optimality for
    layered receivers, not a heuristic assumption.

  \item \textbf{Methodology Banach Floor} (\Cref{thm:banach}): Every
    methodology has a unique attractor $\Sfloor(M)=\sigma\kappa/(1-\kappa)
    >0$ approached geometrically but never attained in finite time.

  \item \textbf{Agent Triple and Uncertainty Principle}
    (\Cref{def:agent,thm:uncertaintyprinciple}): The agent $A=(R,M,G)$
    has floor $\beta\cdot\sigma\kappa/((1-\kappa)\Sigma)$; knowledge and
    action dispersions satisfy a product bound.

  \item \textbf{Bounded Rationality as Theorem}
    (\Cref{thm:boundedrationalitytheorem}): Simon's satisficing is the
    uniquely optimal strategy for a bounded receiver, derived without
    behavioral assumptions.

  \item \textbf{Eleven-fold collapse} (\Cref{thm:elevenfold}): Eleven
    classical prerequisites for point-meaning each individually force
    $\beta=0$, contradicting the Floor Theorem.

  \item \textbf{G\"odelian residue equals private information}
    (\Cref{thm:godelresiduefloor,cor:privateirred}): The irreducible
    private information of any bounded agent equals exactly $\beta$.

  \item \textbf{Classical impossibilities as corollaries}
    (\Cref{sec:impossibilities}): Arrow's impossibility, the
    Grossman--Stiglitz paradox, and the exploration--exploitation
    incompatibility follow as theorems in the present framework.
\end{enumerate}

All thirty-five numerical experiments confirm the quantitative predictions
to relative error $\leq 10^{-14}$.

A companion paper derives market equilibrium conditions from the receiver
primitive: the equilibrium price is a fixed point of a population of
interacting layered receivers, and its existence follows from a collective
version of \Cref{thm:banach}.  All results of that paper are derived
without assumptions beyond those made here, confirming that the receiver
is a sufficient primitive for a complete mathematical theory of economic
exchange.

%------------------------------------------------------------------
%  Bibliography
%------------------------------------------------------------------
\bibliographystyle{plain}
\bibliography{references}

\end{document}